% ============ Chapter 10. I-smoothness (原书 pp.213-245) ============
\chapter{$I$-smoothness}\label{ch:10}

$I$-smoothness is a notion which Grothendieck obtained by reformulating the
theory of simple (non-singular) points in algebraic geometry in terms of an
algebraic `infinitesimal analysis', which makes effective use of nilpotent
elements. The definition looks complicated at first sight, but it has various
alternative formulations, and is a natural and useful notion. In \S\ref{sec:28},
along with the general theory of $I$-smoothness following [G1] we prove the
existence of a coefficient field for a complete local ring of equal
characteristic, relating this to the author's idea of quasi-coefficient field
(Theorem~\ref{thm:28.3}), and discuss Faltings' very simple proof of the
equivalence of $m$-smoothness and geometric regularity for local rings. In
\S\ref{sec:29} we deduce the existence of a coefficient ring for a complete local ring
of unequal characteristic from Theorem~\ref{thm:28.10}, and prove some classical
theorems of Cohen on complete local rings; these results are of decisive
significance for the usefulness of taking completions. \S\ref{sec:30} is something of a
jumble of various theories, but is for the most part occupied with the
so-called Jacobian criterion for regularity. On this subject, we treat the
simple and powerful method obtained by the author's 1972 seminar in the case
of a ring containing a field of characteristic $0$; in the most difficult case of
power series rings in characteristic $p$, the only method currently available
is that of Nagata, and we explain this as simply as possible.

% ============ §28 I-smoothness (原书 pp.213-222) ============
\section{$I$-smoothness}\label{sec:28}\index{I-smooth, section 28}\index{I-unramified, section 28}\index{I-étale, section 28}\index{coefficient field, section 28}\index{quasi-coefficient field, section 28}\index{geometrically regular, section 28}

Let $A$ be a ring, $B$ an $A$-algebra, and $I$ an ideal of $B$; we consider $B$ with the
$I$-adic topology. We say that $B$ is \emph{$I$-smooth} over $A$ if given an $A$-algebra $C$,
an ideal $N$ of $C$ satisfying $N^2 = 0$, and an $A$-algebra homomorphism
$\bar{u}\colon B \to C/N$ which is continuous for the discrete topology of $C/N$ (that is,
such that $\bar{u}(I^\nu) = 0$ for some $\nu$), then there exists a lifting $u\colon B \to C$ of
$\bar{u}$ to $C$.
\[
\begin{tikzcd}
B \arrow[r, "\bar{u}"] & C/N \\
A \arrow[u] \arrow[r] & C \arrow[u]
\end{tikzcd}
\]
If $I = (0)$, that is if no continuity condition is imposed on $\bar{u}$, then this is the
definition of 0-smooth given in \S\ref{sec:25}. Write $f'\colon C \to C/N$ for the natural map;
then from $f'u(I^\nu) = \bar{u}(I^\nu) = 0$ we have $u(I^\nu) \subset N$, hence
$u(I^{2\nu}) \subset N^2 = 0$, so that $u\colon B \to C$ (assuming it exists) is continuous for
the discrete topology of $C$. From this, one sees that if $B$ is $I$-smooth over $A$, and
instead of the condition $N^2 = 0$ we assume that $C$ is an $N$-adically complete
ring, then a continuous homomorphism $\bar{u}\colon B \to C/N$ has a lifting $u\colon B \to C$,
and $u$ is continuous with respect to the $N$-adic topology of $C$; this is because
we can lift $\bar{u}$ successively to $B \to C/N^i$ for $i = 1, 2, \ldots$, and then $u$ is given
by $B \to \varprojlim C/N^i = C$.

We now return to the original assumption $N^2 = 0$; we say that $B$ is
\emph{$I$-unramified} (or $I$-neat) over $A$ if given $C$, $N$ and a continuous
homomorphism $\bar{u}\colon B \to C/N$, there exists at most one lifting of $\bar{u}$ to $C$. If
$B$ is both $I$-smooth and $I$-unramified over $A$, we say that $B$ is \emph{$I$-\'etale}.
These conditions become weaker if we replace $I$ by a larger ideal.

\begin{theorem}[Transitivity]\label{thm:28.1}
Let $A \to B \xrightarrow{g'} B'$ be ring homomorphisms, and suppose that $g'$ is
continuous for the $I$-adic topology of $B$ and the $I'$-adic topology of $B'$; if $B$ is
$I$-smooth over $A$, and $B'$ is $I'$-smooth over $B$ then $B'$ is $I'$-smooth over $A$.
The same thing holds with $I$-unramified in place of $I$-smooth.
\end{theorem}

\begin{proof}
Suppose that $\bar{u}$ is given in the diagram;
\[
\begin{tikzcd}
B' \arrow[r, "\bar{u}"] & C/N \\
B \arrow[u, "g'"] \arrow[ru, "\overline{ug'}"'] & C \arrow[u] \\
A \arrow[u] \arrow[ruu] & 
\end{tikzcd}
\]
then since $\bar{u}g'\colon B \to C/N$ is continuous, by the $I$-smoothness of $B$, there
exists a lifting $w\colon B \to C$. Next by the $I'$-smoothness of $B'$ over $B$, we can
lift $\bar{u}$ to $v\colon B' \to C$. Also if $B$ is $I$-unramified over $A$, and the map $v$ in
the diagram exists, then $w = vg'$ is unique, and if in addition $B'$ is
$I'$-unramified over $B$ then $v$ is unique.
\end{proof}

\begin{theorem}[Base-change]\label{thm:28.2}
Let $A$ be a ring, $B$ and $A'$ two $A$-algebras, and set $B' = B \otimes_A A'$. If $B$ is
$I$-smooth over $A$ then $B'$ is $IB'$-smooth over $A'$. The same thing holds for
$I$-unramified.
\end{theorem}

\begin{proof}
We have the diagram
\[
\begin{tikzcd}
B \arrow[r, "p"] & B' \arrow[r, "\bar{u}"] & C/N \\
A \arrow[u] \arrow[r] & A' \arrow[u] \arrow[ru, ""] & C \arrow[u]
\end{tikzcd}
\]
where $p$, $q$ are the natural homomorphisms. Then if $\bar{u}$ satisfies
$\bar{u}(I^\nu B') = 0$, there is a lifting $v\colon B \to C$ of $\bar{u}p$. Then if we define
$v'\colon B' = B \otimes_A A' \to C$, by $v' = v \otimes_A \mathrm{id}$, this is a lifting of $\bar{u}$ to $C$. For
unramified, this is clear from the fact that $v'$ is uniquely determined by $v$.
\end{proof}

\begin{example*}
Let $k$ be a ring, $(A, m)$ a local ring, $(\hat{A}, \hat{m})$ its completion, and
$k \to A$ a homomorphism. Then
\begin{enumerate}[(i)]
  \item $\hat{A}$ is $\hat{m}$-\'etale over $A$;
  \item $A$ is $m$-smooth (or $m$-unramified) over $k$ $\iff$ $\hat{A}$ is $\hat{m}$-smooth (or
        $\hat{m}$-unramified) over $k$.
\end{enumerate}
We get a proof at once from the fact that $A/m^\nu \cong \hat{A}/\hat{m}^\nu$ for all $\nu$.
\end{example*}

\begin{example*}
Let $A$ be any ring, and set $B = A[X_1, \ldots, X_n]$ and $I = \sum X_iB$; we give $B$
the $I$-adic topology. Then $B$ is $I$-smooth over $A$.
\end{example*}

\begin{remark}
The gap between $I$-smoothness and 0-smoothness has been studied by Tanimoto
[1], [2]. For instance, if $k$ is a field, then $k[X_1, \ldots, X_n]$ is 0-smooth over
$k$ only when $\chark k = p$ and $[k : k^p] < \infty$.
\end{remark}

Let $(A, m, K)$ be a local ring. If $A$ is of characteristic $p$, then also $\chark K = p$;
moreover, if $\chark K = 0$, then $\chark A = 0$, and $A$ contains the rational number
field $\mathbb{Q}$. In either of these cases $A$ is said to be \emph{equicharacteristic}, or a
local ring of \emph{equal characteristic}; this is equivalent to saying that $A$ contains a
field. If $A$ is not of equal characteristic, then either
\[
  \chark A = 0 \text{ and } \chark K = p,
\]
or
\[
  \chark A = p^n \text{ for some } n > 1 \text{ and } \chark K = p.
\]
In this case we say that $A$ is a local ring of \emph{unequal characteristic}.

Let $A$ be an equicharacteristic local ring and let $K'$ be a subfield of $A$. We say
that $K'$ is a \emph{coefficient field} of $A$ if $K'$ maps isomorphically to $K$ under the
natural map $A \to A/m = K$, or equivalently, if $A = K' + m$. Moreover, we say that
$K'$ is a \emph{quasi-coefficient field} of $A$ if $K$ is 0-\'etale over $K'$ (or rather, over
the image of $K'$ in $K$).

\begin{theorem}\label{thm:28.3}
Let $(A, m, K)$ be an equicharacteristic local ring. Then
\begin{enumerate}[(i)]
  \item $A$ has a quasi-coefficient field;
  \item if $A$ is complete, it has a coefficient field;
  \item if the residue field $K$ of $A$ is separable over a subfield $k \subset A$ then
        $A$ has a quasi-coefficient field $K'$ containing $k$;
  \item if $K'$ is a quasi-coefficient field of $A$, then there exists a unique
        coefficient field $K''$ of the completion $\hat{A}$ containing $K'$.
\end{enumerate}
\end{theorem}

\begin{proof}
(iii) Suppose that $B = \{\tau_1, \tau_2, \ldots\}$ is a differential basis of $K/k$, and for
each $\tau_i$, choose an inverse image $\xi_i \in A$. Then by
Theorem~\hyperref[thm:26.8]{26.8}, $\tau_1, \tau_2, \ldots$ are algebraically independent over $k$, so
that the subring $k[\xi_1, \xi_2, \ldots]$ of $A$ meets $m$ in $(0)$, and hence $A$ contains
the field $K' = k(\xi_1, \xi_2, \ldots)$. We identify $K'$ with its image $k(B)$ in $K$, so that
$K$ is clearly 0-\'etale over $K'$, and $K'$ is a quasi-coefficient field, as required.

(i) By assumption $A$ contains a field, so that it contains a perfect field (for
example, the prime subfield). We need only apply (iii) to this.

(iv) In the diagram
\[
\begin{tikzcd}
K \arrow[r, "\cong"] & \hat{A}/\hat{m} \\
K' \arrow[u, "\subset"] \arrow[ru, "\subset"'] & \hat{A} \arrow[u, "p"] \\
& A \arrow[u, "\subset"]
\end{tikzcd}
\]
there exists a unique lifting of the identity map $K \to \hat{A}/\hat{m}$ to $K \to \hat{A}$,
and its image is the required coefficient field.

(ii) follows from (i) and (iv).
\end{proof}

The next lemma will be made more precise in Theorem~\hyperref[thm:28.7]{7}.

\begin{lemma}\label{lem:28.1}
Suppose that $(A, m, K)$ is a Noetherian local ring containing a field $k$. If $A$ is
$m$-smooth over $k$ then $A$ is regular. The converse holds if the residue field $K$
is separable over $k$.
\end{lemma}

\begin{proof}
Take a perfect subfield $k_1 \subset k$; then $k$ is 0-smooth over $k_1$, so that by
transitivity, $A$ is also $m$-smooth over $k_1$, so that we can assume that $k$ is a
perfect field. Also, replacing $A$ by $\hat{A}$, we can assume that $A$ is complete.
Then $A$ has a coefficient field containing $k$; for ease of notation we write $K$
for this, and identify it with the residue field. If $\{x_1, \ldots, x_n\}$ is a minimal
basis of $m$ then as $K$-algebras we have
\[
  A/m^2 \cong K[X_1, \ldots, X_n]/(X_1, \ldots, X_n)^2.
\]
The composite
\[
  A \to A/m^2 \xrightarrow{\cong} K[X_1, \ldots, X_n]/(X)^2 \xrightarrow{\subset} K[X_1, \ldots, X_n]/(X)^2
\]
lifts to $A \to K[[X_1, \ldots, X_n]]$, and by Theorem~\hyperref[thm:8.4]{8.4}, this is surjective. Thus
$\dim A \ge \dim K[[X_1, \ldots, X_n]] = n$, and together with $\embdim A = n$ this
gives the regularity of $A$.

Conversely, if $A$ is regular and $K$ is separable over $k$, then $A$ has a
coefficient field $K$ containing $k$. Let $\{x_1, \ldots, x_n\}$ be a regular system of
parameters of $A$, and define a homomorphism of $K$-algebras
$\psi\colon K[[X_1, \ldots, X_n]] \to \hat{A}$ by $\psi(X_i) = x_i$; then once more by
Theorem~\hyperref[thm:8.4]{8.4}, $\psi$ is surjective, and comparing dimensions, we see that
\[
  K[[X_1, \ldots, X_n]] = \hat{A}.
\]
Therefore $\hat{A}$ is $\hat{m}$-smooth over $K$, and since $K$ is 0-smooth over $k$, we see
that $\hat{A}$ is $\hat{m}$-smooth over $k$, and therefore $A$ is $m$-smooth over $k$.
\end{proof}

Let $k \to A \xrightarrow{g} B$ be ring homomorphisms, and let $I$ be an ideal of $B$; we
consider $B$ in the $I$-adic topology. We say that $B$ is \emph{$I$-smooth over $A$ relative
to $k$} if the following condition holds: for any $A$-algebra $C$, and an ideal $N$ of
$C$ such that $N^2 = 0$, given an $A$-algebra homomorphism $\bar{u}\colon B \to C/N$
satisfying $\bar{u}(I^\nu) = 0$ for sufficiently large $\nu$, if $\bar{u}$ has a lifting
$v'\colon B \to C$ as a $k$-algebra homomorphism, it also has a lifting $v\colon B \to C$ as an
$A$-algebra homomorphism.

\begin{theorem}\label{thm:28.4}
Let $k \to A \xrightarrow{g} B$ and $I \subset B$ be as above; then the following conditions
are equivalent:
\begin{enumerate}[(1)]
  \item $B$ is $I$-smooth over $A$ relative to $k$;
  \item if $N$ is a $B$-module such that $I^\nu N = 0$ for sufficiently large $\nu$, then
        $\Der_k(B, N) \to \Der_k(A, N)$ is surjective;
  \item for every $\nu > 0$, the map $\Omega_{B/k} \otimes_B (B/I^\nu) \to
        \Omega_{A/k} \otimes_A (B/I^\nu)$ has a left inverse (that is, it maps injectively onto
        a direct summand).
\end{enumerate}
\end{theorem}

\begin{proof}
(1) $\Rightarrow$ (2) If $I^\nu N = 0$, set $C = (B/I^\nu) \ltimes N$, and let
$\bar{u}\colon B \to B/I^\nu = C/N$ be the natural map. Given $D \in \Der_k(A, N)$, define
$\lambda\colon A \to C$ by $\lambda(a) = (ug(a), D(a))$. If we consider $C$ as an $A$-algebra via
$\lambda$, then $b \mapsto (u(b), 0) \in C$ is a $k$-algebra homomorphism from $B$ to $C$
lifting $\bar{u}$, so that by assumption there exists a lifting $v'\colon B \to C$ of $\bar{u}$ as an
$A$-algebra homomorphism; then writing $v'$ in the form
\[
  v'(b) = (u(b), D'(b)), \quad \text{with } D' \in \Der_k(B, N),
\]
we have $v'g = \lambda$, so that $D'g = D$.

(2) $\Rightarrow$ (1) Suppose given a commutative diagram
\[
\begin{tikzcd}
B \arrow[r, "\bar{u}"] & C/N \\
A \arrow[u, "g"] \arrow[ru, ""] & C \arrow[u, "j"]
\end{tikzcd}
\]
with $j$ the natural map, and $\bar{u}(I^\nu) = 0$; if $v\colon B \to C$ is a $k$-algebra
homomorphism satisfying $jv = \bar{u}$ and $vg = \lambda$, then setting $D = \lambda - vg$, we
can view $D$ as an element of $\Der_k(A, N)$. By assumption there exists
$D' \in \Der_k(B, N)$ such that $D = D'g$. Using this, we set $v' = v + D'$; then
\[
  v'g = vg + D'g = \lambda - D + D = \lambda, \quad \text{and} \quad jv' = \bar{u}.
\]

(2) $\Leftrightarrow$ (3) comes from observing the general fact that for a ring $R$, a map
$\varphi\colon M \to M'$ of $R$-modules has a left inverse if and only if for every $R$-module
$N$ the induced map
\[
  \Hom_R(M', N) \to \Hom_R(M, N)
\]
is surjective.
\end{proof}

\begin{theorem}\label{thm:28.5}
Let $A$ be a ring, $B$ an $A$-algebra, and $I$ an ideal of $B$; set $\bar{B} = B/I$, and
assume that $B$ is $I$-smooth over $A$. Then $\Omega_{B/k} \otimes_B \bar{B}$ is projective as
a $\bar{B}$-module.
\end{theorem}

\begin{proof}
It is enough to show that for an exact sequence $L \to M \to 0$ of $\bar{B}$-modules, the
sequence
\[
  \Hom_{\bar{B}}(\Omega_{B/k} \otimes_B \bar{B}, L) \to \Hom_{\bar{B}}(\Omega_{B/k} \otimes_B \bar{B}, M) \to 0
\]
is exact, that is that
\[
  \Der_k(B, L) \to \Der_k(B, M) \to 0
\]
is exact. Set $C = \bar{B} \ltimes L$ and $N = \Ker \varphi$. If we view both $L$ and $N$ as
ideals of $C$, we have $L^2 = N^2 = 0$ and $C/N \cong \bar{B} \ltimes M$. Now for any
$D \in \Der_k(B, M)$ we have an $A$-algebra homomorphism $B \to C/N$ given by
\[
  b \mapsto (b, D(b)) \in \bar{B} \ltimes M,
\]
and lifting this to $B \to C$ is equivalent to lifting $D$ to an element of
$\Der_k(B, L)$.
\end{proof}

\begin{lemma}\label{lem:28.2}
Let $B$ be a ring and $I$ an ideal of $B$, and let $u\colon L \to M$ be a map of
$B$-modules; assume that $M$ is projective. Suppose also that one of the
following two conditions hold:
\begin{enumerate}[(a)]
  \item $I$ is nilpotent;
or ($\beta$) $L$ is a finite $B$-module and $I \subset \rad(B)$.
\end{enumerate}
Then
\[
  u \text{ has a left inverse} \iff \bar{u}\colon L/IL \to M/IM \text{ has a left inverse}.
\]
\end{lemma}

\begin{proof}
($\Leftarrow$) is trivial. To prove ($\Rightarrow$), suppose that $\bar{v}\colon M/IM \to L/IL$ is a
left inverse of $\bar{u}$. Since $M$ is projective, there is a map $w\colon M \to L$ such
that the diagram
\[
\begin{tikzcd}
M \arrow[r, "w"] \arrow[d] & L \arrow[d] \\
M/IM \arrow[r, "\bar{w}"'] & L/IL
\end{tikzcd}
\]
commutes. Set $\tau = uw$. Then $\tau$ induces the identity on $L/IL$, so that
$L = \tau(L) + IL$, so that by NAK, $L = \tau(L)$. Hence if $L$ is a finite $B$-module,
then by Theorem~\ref{thm:2.4}, $\tau$ is also injective. Furthermore, if $I^\nu = 0$ we do
the following: if $x \in \Ker \tau$ then $0 = \bar{\tau}(\bar{x}) = \bar{x} \mod IL$, so that $x \in IL$, and
we can write $x = \sum a_i y_i$ with $a_i \in I$ and $y_i \in L$. Then
\[
  0 = \tau(x) = \sum a_i \tau(y_i) \equiv \sum a_iy_i = x \mod I^2L,
\]
so that $x \in I^2L$, and proceeding in the same way we arrive at
$x \in I^\nu L = 0$. Hence also in this case $\tau$ is an automorphism of $L$, and
$\tau^{-1}u$ is the required left inverse of $u$.
\end{proof}

\begin{theorem}\label{thm:28.6}
Let $k \to A \to B$ be ring homomorphisms, $I$ an ideal of $B$, and suppose that $B$ is
$I$-smooth over $k$. Set $B_1 = B/I$. Then the following conditions are equivalent:
\begin{enumerate}[(1)]
  \item $B$ is $I$-smooth over $A$;
  \item $\Omega_{A/k} \otimes_A B_1 \to \Omega_{B/k} \otimes_B B_1$ has a $B_1$-linear left inverse.
\end{enumerate}
\end{theorem}

\begin{proof}
(1) $\Rightarrow$ (2) is contained in Theorem~\hyperref[thm:28.4]{4}. Conversely, suppose that (2)
holds. For any $\nu > 0$, set $B_\nu = B/I^\nu$; then since $I$-smoothness and
$I^\nu$-smoothness are the same, by Theorem~\hyperref[thm:28.5]{5}, $\Omega_{B/k} \otimes_B B_\nu$ is a
projective $B_\nu$-module. Now set $I_1 = I/I^\nu$; then $B_\nu/I_1 = B_1$, and
$I_1^2 = 0$, so that applying Lemma~\hyperref[lem:28.2]{2}, we see that
$\Omega_{A/k} \otimes_A B_\nu \to \Omega_{B/k} \otimes_B B_\nu$ has a left inverse. By
Theorem~\hyperref[thm:28.4]{4}, $B$ is $I$-smooth over $A$ relative to $k$, but since it is also
$I$-smooth over $k$, it is also $I$-smooth over $A$.
\end{proof}

\begin{corollary*}
Let $(A, m, K)$ be a regular local ring containing a field $k$; then
\[
  A \text{ is } m\text{-smooth over } k \iff \Omega_k \otimes_k K \to \Omega_A \otimes_A K \text{ is injective}.
\]
\end{corollary*}

\begin{proof}
Let $k_1 \subset k$ be the prime subfield. Then by Lemma~\hyperref[lem:28.1]{1}, $A$ is $m$-smooth
over $k_1$, so that we need only apply the theorem to $k_1 \to k \to A$.
\end{proof}

Let $A$ be a Noetherian local ring, and $k \subset A$ a subfield. We say that $A$ is
\emph{geometrically regular} over $k$ if $A \otimes_k k'$ is a regular ring for every finite
extension field $k'$ of $k$.

\begin{theorem}\label{thm:28.7}
Let $(A, m, K)$ be a Noetherian local ring, and $k \subset A$ a subfield; then
\[
  A \text{ is } m\text{-smooth over } k \iff A \text{ is geometrically regular over } k.
\]
\end{theorem}

\begin{proof}
($\Rightarrow$) Let $k'$ be a finite extension field of $k$. Then $A \otimes_k k' = A'$ is
$mA'$-smooth over $k'$ by base-change. Let $n$ be any maximal ideal of $A'$; then $A'$ is
a finite $A$-module, hence integral over $A$, so that $n \supset mA'$. Thus if we set
$A'' = A'_n$ and $m'' = nA''$ then $A' \to A''$ is continuous for the $mA'$-adic topology of
$A'$ and the $m''$-adic topology of $A''$, but as a localisation $A''$ is 0-\'etale over $A'$,
so that by Theorem~\hyperref[thm:28.1]{1}, $A''$ is $m''$-smooth over $k'$, and hence by
Lemma~\hyperref[lem:28.1]{1}, $A'' = A'_n$ is a regular local ring. This is what was required
to prove.

($\Leftarrow$) According to Lemma~\hyperref[lem:28.1]{1}, there is only a problem if $k$ is of
characteristic $p$. The proof we now give was discovered in 1977 by G.~Faltings
[1] while he was still a student.

By the corollary of the previous theorem, we need only prove that
$\Omega_k \otimes_k K \to \Omega_A \otimes_A K$ is injective. For this, we let
$x_1, \ldots, x_r \in k$ be $p$-independent elements, and prove that
$dx_1, \ldots, dx_r \in \Omega_A \otimes_A K$ are linearly independent over $K$. Write $\alpha_i$ for $p$-th
roots of the $x_i$, and set $k' = k(\alpha_1, \ldots, \alpha_r)$. Then
\[
  B = A \otimes_k k' = A[T_1, \ldots, T_r]/(T_1^p - x_1, \ldots, T_r^p - x_r)
\]
is a Noetherian local ring; write $n$ for its maximal ideal, and $L$ for the residue
field $L = B/n$. By Theorem~\hyperref[thm:25.2]{25.2} the sequence
\[
  0 \to n/n^2 \to \Omega_{B/k} \otimes_B L \to \Omega_{L/k} \to 0
\]
is exact. Similarly,
\[
  0 \to m/m^2 \to \Omega_{A/k} \otimes_A K \to \Omega_{K/k} \to 0
\]
is exact. Now consider the commutative diagram
\[
\begin{tikzcd}
0 \arrow[r] & n/n^2 \arrow[r] \arrow[d, "\varphi_1"'] & \Omega_{B/k} \otimes_B L \arrow[r] \arrow[d, "\varphi_2"'] & \Omega_{L/k} \arrow[r] \arrow[d, "\varphi_3"'] & 0 \\
0 \arrow[r] & (m/m^2) \otimes_K L \arrow[r] & \Omega_{A/k} \otimes_A L \arrow[r] & \Omega_{K/k} \otimes_K L \arrow[r] & 0
\end{tikzcd}
\]
where $\varphi_1$, $\varphi_2$ and $\varphi_3$ are the natural maps. Then by the snake lemma, we get a
long exact sequence of $L$-modules
\[
  0 \to \Ker \varphi_1 \to \Ker \varphi_2 \to \Ker \varphi_3
      \to \Coker \varphi_1 \to \Coker \varphi_2 \to \Coker \varphi_3 \to 0.
\]
By assumption $A$ and $B$ are regular local rings of the same dimension, so that
$\rank m/m^2 = \dim A = \rank n/n^2$, so $\rank \Ker \varphi_1$ and $\rank \Coker \varphi_1$ are
finite and equal; moreover, since $L$ is a finite extension field of $K$, both
$\rank \Ker \varphi_3$ and $\rank \Coker \varphi_3$ are finite and equal (the Cartier equality).
Therefore by the above exact sequence, we get
\[
  \rank \Ker \varphi_2 = \rank \Coker \varphi_2.
\]
However, $\Coker \varphi_2 = \Omega_{B/A} \otimes_B L$, and by Theorem~\hyperref[thm:25.2]{25.2},
\[
  \Omega_{B/A} = BdT_1 + \cdots + BdT_r,
\]
so that both of $\Ker \varphi_2$ and $\Coker \varphi_2$ have rank equal to $r$. Now if we set
$J = (T_1^p - x_1, \ldots, T_r^p - x_r)$ we have an exact sequence
\[
  J/J^2 \to \Omega_{A[T_1, \ldots, T_r]/k} \otimes_A B = \Omega_{A/k} \otimes_A B \oplus B dT_i \to \Omega_{B/k} \to 0,
\]
and since this remains exact after performing $\otimes_A L$, we see that $\Ker \varphi_2$ is
generated by $dx_1, \ldots, dx_r$. Therefore $dx_1, \ldots, dx_r \in \Omega_A \otimes_A L$ are
linearly independent over $L$, so that they must also be linearly independent
over $K$ as elements of $\Omega_A \otimes_A K$. This is what we had to prove.
\end{proof}

Let $B$ be an $A$-algebra, $I$ an ideal of $B$, and consider $B$ with the $I$-adic
topology. Let $N$ be a $B$-module such that $I^\nu N = 0$ for some $\nu > 0$; in what
follows we will say that a $B$-module with this property is \emph{discrete}. An
$A$-bilinear map $f\colon B \times B \to N$ will be called a \emph{continuous symmetric
2-cocycle} if it satisfies the three conditions
\begin{enumerate}[(a)]
  \item $xf(y, z) - f(xy, z) + f(x, yz) - f(x, y)z = 0$ for all $x, y, z \in B$,
  \item[($\beta$)] $f(x, y) = f(y, x)$,
  \item[($\gamma$)] there exists $\rho \ge \nu$ such that $f(x, y) = 0$ if either $x \in I^\rho$ or $y \in I^\rho$.
\end{enumerate}
If this holds, we set $f(1, 1) = \varepsilon$; then substituting $y = z = 1$ in (a) gives
$x\varepsilon = f(x, 1)$.

Define a product on the $A$-module $C = (B/I^\nu) \oplus N$ by
\[
  (x, v)(y, w) = (xy, -f(x, y) + xw + yv)
\]
for $x, y \in B$; then $C$ is a commutative ring with unit $(1, \varepsilon)$, and $N$ is an
ideal of $C$ satisfying $N^2 = 0$. If we define a map $A \to C$ by $a \mapsto (\bar{a}, a\varepsilon)$
then this is a ring homomorphism, and the diagram
\[
\begin{tikzcd}
B \arrow[r] & B/I^\nu = C/N \\
A \arrow[u] \arrow[r] & C \arrow[u]
\end{tikzcd}
\]
is commutative; then a necessary and sufficient condition for $\bar{u}$ to have a lifting
to $B \to C$ is that there should exist an $A$-linear map $g\colon B \to N$ such that
\[
  (\delta) \quad f(x, y) = xg(y) - g(xy) + g(x)y \quad \text{for all } x, y \in B.
\]
For if $g$ exists, then defining $v\colon B \to C$ by $v(x) = (\bar{x}, g(x))$ we find that $v$
is a lifting of $\bar{u}$, and conversely, if there is a lifting $v$ of $\bar{u}$ one checks easily
that one can find a $g$ as above.

We say that the 2-cocycle $f$ \emph{splits} if there exists $g$ satisfying ($\delta$). For any
$A$-linear map $g\colon B \to N$, we write $\delta g$ for the bilinear map
$B \times B \to N$ given by the right-hand side of ($\delta$); this satisfies (a) and ($\beta$),
and if $g$ is continuous (that is, $\exists \rho$ such that $g(I^\rho) = 0$), then it also
satisfies ($\gamma$).

\begin{theorem}\label{thm:28.8}
Let $A$ be a ring and $B$ an $A$-algebra with an $I$-adic topology.
\begin{enumerate}[(i)]
  \item If $B$ is $I$-smooth over $A$ then every continuous symmetric 2-cocycle
        $f\colon B \times B \to N$ with values in a discrete $B$-module $N$ splits.
  \item If $B/I^n$ is projective as an $A$-module for infinitely many $n$, and if every
        continuous symmetric 2-cocycle with values in a discrete $B$-module splits,
        then $B$ is $I$-smooth over $A$.
\end{enumerate}
\end{theorem}

\begin{proof}
(i) is what we have just said.

(ii) Suppose that we are given a commutative diagram
\[
\begin{tikzcd}
B \arrow[r, "\bar{u}"] & C/N \\
A \arrow[u] \arrow[r] & C \arrow[u]
\end{tikzcd}
\]
with $N^2 = 0$ and $\bar{u}(I^\nu) = 0$; then in view of $N^2 = 0$, the $C$-module $N$ can be
viewed as a $C/N$-module, and by means of $\bar{u}$ as a $B$-module; but then
$I^\nu N = 0$, so that $N$ is a discrete $B$-module. Take an integer $n > \nu$ such that
$B/I^n$ is projective as an $A$-module. Then $\bar{u}$ can be lifted as a map of $A$-modules
to $\lambda\colon B \to C$ such that $\lambda(I^n) = 0$. For $x, y \in B$ we set
\[
  f(x, y) = \lambda(xy) - \lambda(x)\lambda(y),
\]
and since $\lambda$ is a ring homomorphism modulo $N$, we have $f(x, y) \in N$. Now for
$\xi \in N$ and $x \in B$, by definition we have $\lambda(x) \cdot \xi = x \cdot \xi$ (both sides are
products evaluated in $C$), and using this one computes the left-hand side of (a) to
be zero. The symmetry ($\beta$) is obvious. Also $\lambda(I^n) = 0$, so that we also get
($\gamma$). Thus $f$ is a continuous symmetric 2-cocycle, and hence by assumption
there is an $A$-linear map $g\colon B \to N$ satisfying
\[
  f(x, y) = xg(y) - g(xy) + g(x)y.
\]
Now if we set $v = \lambda + g$, we have
\begin{align*}
  v(xy) &= \lambda(xy) + g(xy) \\
        &= \lambda(x)\lambda(y) + f(x, y) + g(xy) \\
        &= \lambda(x)\lambda(y) + xg(y) + g(x)y \\
        &= (\lambda(x) + g(x))(\lambda(y) + g(y)) = v(x)v(y),
\end{align*}
so that $v$ is an $A$-algebra homomorphism, and is a lifting of $\bar{u}$.
\end{proof}

\begin{theorem}[([G1], 19.7.1)]\label{thm:28.9}
Let $(A, m, k)$ and $(B, n, k')$ be Noetherian local rings, and $\varphi\colon A \to B$ a
local homomorphism; set $B_0 = B \otimes_A k = B/mB$ and $n_0 = n/mB$. Then the
following conditions are equivalent:
\begin{enumerate}[(1)]
  \item $B$ is $n$-smooth over $A$,
  \item $B$ is flat over $A$ and $B_0$ is $n_0$-smooth over $k$.
\end{enumerate}
\end{theorem}

\begin{proof}
This is an extremely important theorem, but the proof is long and difficult, and
we refer to [G1] for it. We content ourselves with proving the following
analogous theorem, which is all that we will need in what follows.
\end{proof}

\begin{theorem}\label{thm:28.10}
Let $(A, m, k)$ be a local ring, and $B$ a flat $A$-algebra; suppose that
$B_0 = B \otimes_A k$ is 0-smooth over $k$. Then $B$ is $mB$-smooth over $A$.
\end{theorem}

\begin{proof}
As one sees from the definition of $mB$-smoothness, it is enough to prove that
$B/m^\nu B$ is 0-smooth over $A/m^\nu$ for every $\nu > 0$. Since $B/m^\nu B$ is flat over
$A/m^\nu$, we can assume that $m$ is nilpotent. Then a flat $A$-module is free (by
Theorem~\ref{thm:7.10}), so that $B$ is a projective $A$-module, and hence by
Theorem~\hyperref[thm:28.8]{8}, we need only show that every symmetric 2-cocycle
$f\colon B \times B \to N$ with values in a $B$-module $N$ splits. First of all, in the case
that $N$ satisfies $mN = 0$, then since $f$ is $A$-bilinear $f$ is essentially a
2-cocycle over $B_0$, that is there is a map $f_0\colon B_0 \times B_0 \to N$ such that
$f(x, y) = f_0(\bar{x}, \bar{y})$.

Now $B_0$ is 0-smooth over $k$, so that by Theorem~\hyperref[thm:28.8]{8}, $f_0$ splits, that is
there is a map $g_0\colon B_0 \to N$ such that $f_0 = \delta g_0$. Thus setting
$g(x) = g_0(\bar{x})$ we get
\[
  f = \delta g.
\]
In the general case, write $\varphi$ for the natural map $N \to N/mN$, and consider
$\varphi \circ f$; then this splits, that is there exists $\bar{g}\colon B \to N/mN$ such that
$\varphi f = \delta \bar{g}$.

Now since $B$ is projective over $A$, we can lift $\bar{g}$ to an $A$-linear map
$g\colon B \to N$, and then $f - \delta g$ is a 2-cocycle with values in $mN$. Doing the same
thing once more, we find $h\colon B \to mN$ such that $f - \delta(g + h)$ is a 2-cocycle with
values in $m^2N$. Proceeding in the same way, since $m$ is nilpotent we finally see
that $f$ splits.
\end{proof}

% ============ 习题 ============
\section*{Exercises to \S\ref{sec:28}}
\emph{Prove the following propositions.}

\begin{exercise}\label{ex:28.1-ex}
Theorem~\hyperref[thm:28.10]{28.10} also holds on replacing smooth by unramified or \'etale.
\end{exercise}

\begin{exercise}\label{ex:28.2-ex}
Let $k$ be a non-perfect field of characteristic $p$, and $a \in k - k^p$; set
$A = k[X]_{(X^p - a)}$. Then the residue field $k(a^{1/p})$ of $A$ is inseparable over
$k$. This ring $A$ does not have a coefficient field containing $k$, but is
0-smooth over $k$.
\end{exercise}

% ============ §29 The structure theorems for complete local rings (原书 pp.223-230) ============
\section{The structure theorems for complete local rings}\label{sec:29}\index{coefficient ring, section 29}\index{p-ring, section 29}\index{quasi-coefficient ring, section 29}\index{Eisenstein extension, section 29}\index{ramified local ring, section 29}\index{unramified local ring, section 29}

By Theorem~\hyperref[thm:28.3]{28.3}, a complete local ring of equal characteristic $A$ has a
coefficient field. If $K$ is a coefficient field of $A$, and $x_1, \ldots, x_n$ are generators
of the maximal ideal, then any element of $A$ can be expanded as a power series
in $x_1, \ldots, x_n$ with coefficients in $K$, and therefore $A$ is a quotient of the
regular local ring $K[[X_1, \ldots, X_n]]$. We now want to extend all of this to the
case of unequal characteristic.

We say that a DVR of characteristic $0$ is a \emph{$p$-ring}\index{p-ring} if its maximal ideal is
generated by the prime number $p$. If $K$ is a given field of characteristic $p$, then
there exists a $p$-ring having $K$ as its residue field. This follows by applying the
next theorem to $A = \mathbb{Z}_{(p)}$.

\begin{theorem}\label{thm:29.1}
Let $(A, tA, k)$ be a DVR and $K$ an extension field of $k$; then there exists a DVR
$(B, tB, K)$ containing $A$.
\end{theorem}

\begin{proof}
Let $\{x_\lambda\}_{\lambda \in \Lambda}$ be a transcendence basis of $K$ over $k$, and set
$k_1 = k(\{x_\lambda\})$. We take indeterminates $\{X_\lambda\}_{\lambda \in \Lambda}$ over $A$ in
bijection with the $\{x_\lambda\}$, and set $A[\{X_\lambda\}] = A'$ and $A_1 = (A')_{tA'}$. Now
$A'$ is a free $A$-module, and hence separated for the $t$-adic topology; therefore,
so is $A_1$, and $A_1$ is a DVR with $A_1/tA_1 \cong k_1$. Hence, replacing $A$ and $k$ by $A_1$
and $k_1$, we can assume that $K$ is algebraic over $k$. Let $L$ be an algebraic
closure of the field of fractions of $A$, and let $\mathcal{F}$ be the set of all pairs
$(B, \varphi)$, where $B$ is an intermediate ring $A \subset B \subset L$ and
$\varphi\colon B \to K$ is an $A$-algebra homomorphism, satisfying the conditions
\[
  (*) \quad B \text{ is a DVR and } \Ker \varphi = \rad(B) \text{ is generated by } t.
\]
We introduce an order on $\mathcal{F}$ by defining $(B, \varphi) \le (C, \psi)$ if $B \subset C$ and
$\psi|_B = \varphi$. Let us show that $\mathcal{F}$ has a maximal element. Suppose that
$\{(B_i, \varphi_i)\}$ is a totally ordered subset of $\mathcal{F}$, and set
\[
  B_0 = \bigcup_i B_i, \quad \varphi_0 = \bigcup_i \varphi_i;
\]
then one sees easily that $B_0$ is a local ring with maximal ideal $tB_0$. If
$0 \ne x \in B_0$, then $x \in B_i$ for some $i$, and since $B_i$ is a DVR we can write
$x = t^\nu u$ with $u$ a unit of $B_i$ and some $\nu$. From this we get
$x \notin t^{\nu+1}B_0$, so that $B_0$ is $t$-adically separated. Hence $B_0$ is a DVR, and if
$\varphi_0\colon B_0 \to K$ is defined to be equal to $\varphi_i$ on $B_i$ then
$(B_0, \varphi_0) \in \mathcal{F}$. Hence by Zorn's lemma $\mathcal{F}$ has a maximal element;
suppose that $(B, \varphi)$ is one. Then if $\varphi(B) \ne K$, take $\alpha \in K - \varphi(B)$,
let $\bar{f}(X)$ be the minimal polynomial of $\alpha$ over $\varphi(B)$, and take a manic
$f(X) \in B[X]$ which is an inverse image of $\bar{f}(X)$. Then $f(X)$ is irreducible in
$B[X]$, and hence (by Ex.~\ref{ex:9.6-ex}) also irreducible over the field of fractions of
$B$. Let $\xi$ be a root of $f(X)$ in $L$, and set $B' = B[\xi]$; then
$B' = B[X]/(f)$. Therefore
\[
  B'/tB' = B[X]/(t, f) = \varphi(B)[X]/(\bar{f}) = \varphi(B)(\alpha)
\]
is a field; since $B'$ is integral over $B$, every maximal ideal of $B'$ contains
$tB'$, so that $B'$ is a local integral domain with maximal ideal $tB'$, and is
Noetherian because it is finite over $B$, therefore a DVR. This contradicts the
maximality of $B$, so that we must have $\varphi(B) = K$.
\end{proof}

\begin{remark}
Since $B$ is an integral domain containing $A$, it is flat over $A$ by
Ex.~\ref{ex:10.2-ex}. In EGA $0_{IV}$ (10.3.1), the following more general fact is proved:
let $(A, m, k)$ be a Noetherian local ring, and $K$ an extension field of $k$, then
there exists a Noetherian ring $B$ containing $A$ satisfying the three conditions
(1) $\rad(B) = mB$, (2) $B/mB$ is isomorphic over $k$ to $K$, and (3) $B$ is flat over $A$.
\end{remark}

\begin{theorem}\label{thm:29.2}
Let $(A, m, K)$ be a complete local ring, $(R, pR, k)$ a $p$-ring, and
$\varphi_0\colon k \to K$ a field homomorphism; then there exists a local homomorphism
$\varphi\colon R \to A$ which induces $\varphi_0$ on the residue fields.
\end{theorem}

\begin{proof}
Set $S = \mathbb{Z}_{(p)}$, and let $k_0 \subset k$ be the prime subfield. Since
$\varphi_0(k_0) \subset K$, the prime number $p$, viewed as an element of $A$, belongs to
$m$. Hence the standard homomorphism $\mathbb{Z} \to A$ extends to a local
homomorphism $S \to A$. Now $R \otimes_S k_0 = R/pR = k$ is a separable extension of
$k_0$, and hence 0-smooth over $k_0$; also $R$ is a torsion-free $S$-module, hence
flat over $S$, so that by Theorem~\hyperref[thm:28.10]{28.10}, $R$ is $pR$-smooth over $S$.
\[
\begin{tikzcd}
k_0 \arrow[r, "\subset"] \arrow[d, "\subset"'] & k \arrow[r, "\varphi_0"] & K = A/m \\
S \arrow[r, "\subset"] & R \arrow[ru, "\varphi"', dashed] & A \arrow[u, "p"] \\
& & \mathbb{Z} \arrow[u, "\subset"] \arrow[uu, "\subset"]
\end{tikzcd}
\]
Hence, as we discussed at the beginning of \S\ref{sec:28}, we can lift
$\varphi_0\colon k \to A/m = K$ successively to $R \to A/m^\nu$, and using the fact that
$A = \varprojlim A/m^\nu$, we get $\varphi\colon R \to A$ making the left-hand diagram
commute.
\end{proof}

\begin{theorem}\label{thm:29.3}
The complete $p$-ring is uniquely determined up to isomorphism by its residue
field.
\end{theorem}

\begin{proof}
Suppose that $R$ and $R'$ are both complete $p$-rings with residue field $k$; then by
the theorem there exists a local homomorphism $\varphi\colon R \to R'$ which induces
the identity map on the residue field. We have $R' = \varphi(R) + pR'$, and of course
$\varphi(p) = p$, so that by the completeness of $R$ we see that $\varphi$ is
surjective, and is also injective, since $p^\nu R$ is not contained in $\Ker \varphi$
for any $\nu$. Therefore $R \cong R'$.
\end{proof}

Let $(A, m, k)$ be a complete local ring of unequal characteristic, and let
$p = \chark k$. We say that a subring $A_0 \subset A$ is a \emph{coefficient ring}\index{coefficient ring}
of $A$ if $A_0$ is a complete Noetherian local ring with maximal ideal $pA_0$, and
\[
  A = A_0 + m, \quad \text{that is, } k = A/m \cong A_0/pA_0.
\]
By Theorem~\hyperref[thm:29.1]{1} applied to the residue field $k$ of $A$, there exists a $p$-ring
$S$ such that $S/pS = k$; write $R$ for the completion of $S$, so that $R$ is a
complete $p$-ring with residue field $k$. By Theorem~\hyperref[thm:29.2]{2}, there exists a local
homomorphism $\varphi\colon R \to A$ inducing an isomorphism on the residue fields. If
we set $\varphi(R) = A_0$, then this is clearly a coefficient ring of $A$. If $A$ has
characteristic $0$ then $\varphi$ is injective and $A_0 \cong R$. If $A$ has
characteristic $p^\nu$ then $A_0 \cong R/p^\nu R$. We summarise the above discussion in
the following theorem.

\begin{theorem}\label{thm:29.4}
If $(A, m, k)$ is a complete local ring and $p = \chark k$ then $A$ has a coefficient
ring $A_0$. If $A$ has characteristic $0$ then $A_0$ is a complete DVR, and if $A$ has
characteristic $p^\nu$ then $A_0$ is an Artinian local ring.
\end{theorem}

In what follows, in order to include the case of equal characteristic in our
discussion, we also consider coefficient fields as being coefficient rings.
From the previous theorem and Theorem~\hyperref[thm:28.3]{28.3}, we get the following
important result.

\begin{theorem}\label{thm:29.5}
\begin{enumerate}[(i)]
  \item If $(A, m)$ is a complete local ring and $m$ is finitely generated, then $A$
        is Noetherian.
  \item A Noetherian complete local ring is a quotient of a regular local ring;
        in particular it is universally catenary.
  \item If $A$ is a Noetherian complete local ring (and in the case of unequal
        characteristic, $A$ is an integral domain), then there exists a subring
        $A' \subset A$ with the following properties: $A'$ is a complete regular local
        ring with the same residue field as $A$, and $A$ is finitely generated as an
        $A'$-module.
\end{enumerate}
\end{theorem}

\begin{proof}
We choose a coefficient ring $A_0$ of $A$. If $m = (x_1, \ldots, x_n)$ then every element
of $A$ can be expanded as a power series in $x_1, \ldots, x_n$ with coefficients in
$A_0$, so that $A$ is a quotient of $A_0[[X_1, \ldots, X_n]]$ and hence Noetherian. Now
$A_0$ is a quotient of a $p$-ring $R$, so that $A$ is a quotient of
$R[[X_1, \ldots, X_n]]$, which is a regular local ring, hence a CM ring, and therefore
according to Theorem~\hyperref[thm:17.9]{17.9}, $A$ is universally catenary. To prove (iii), set
$\dim A = n$, and in the case of equal characteristic let $\{y_1, \ldots, y_n\}$ be any
system of parameters of $A$; if $A$ is an integral domain of characteristic $0$ and
$\chark k = p$, we can choose a system of parameters $\{y_1 = p, y_2, \ldots, y_n\}$ of $A$
starting with $p$. In either case $R = A_0$, so that we set $A' = R[y]$. Then $A'$ is the
image of $\varphi\colon R[Y] \to A$, where $R[Y]$ is the regular local ring
\[
  R[Y] = R[[Y_1, \ldots, Y_n]], \quad \text{or } R[[Y_2, \ldots, Y_n]] \text{ if } y_1 = p,
\]
and $\varphi$ is the $R$-algebra homomorphism defined by $\varphi(Y_i) = y_i$. Set
$m' = \sum y_iA'$. Since $A/m = A'/m'$, every $A$-module of finite length has the same
length when viewed as an $A'$-module. In particular $A/m'A$ is a finite module over
$A'/m'$, and $A$ is $m'$-adically separated, so that by Theorem~\hyperref[thm:8.4]{8.4} $A$ is a finite
$A'$-module. Therefore, we have
\[
  \dim A' = \dim A = n.
\]
$R[Y]$ is an $n$-dimensional integral domain, and if $\Ker \varphi \ne 0$ we would have
$\dim A' < n$, which is a contradiction. Therefore $\varphi$ is injective and
$A' \cong R[Y]$.
\end{proof}

\begin{remark}
In the case of unequal characteristic when $A$ is not an integral domain, (iii)
can fail to hold. If $A$ is of characteristic $p^m$ with $m > 1$, then every subring
of $A$ has the same characteristic $p^m$, so cannot be regular. Even if $A$ has
characteristic $0$, the following is a counter-example: let $R$ be a complete
$p$-ring and $A = R[X]/(pX)$; then $A$ is a complete one-dimensional Noetherian
local ring, but if a subring $A'$ as in (iii) were to exist, $A'$ would be a
one-dimensional regular local ring, hence a DVR, and since $A'$ has characteristic
$0$ and its residue field characteristic $p$, $A'/pA'$ would be an Artinian ring,
and hence also $A/pA$ would be Artinian. But $A/pA \cong k[X]$ is
one-dimensional, and this is a contradiction.

The proof of Theorem~\hyperref[thm:29.5]{5} shows that it is sufficient to assume that $p$ is not
in any minimal prime ideal of $A$.
\end{remark}

Corresponding to the definition of quasi-coefficient field of an
equicharacteristic local ring, let us define \emph{quasi-coefficient rings} in the case of
unequal characteristic. Let $(A, m, K)$ be a possibly non-complete local ring, and
suppose $\chark K = p$. A subring $S \subset A$ is said to be a quasi-coefficient ring if
it satisfies the following two conditions:
\begin{enumerate}[(1)]
  \item $S$ is a Noetherian local ring with maximal ideal $pS$;
  \item $K = A/m$ is 0-\'etale over $S/pS$.
\end{enumerate}
In view of (1) if $A$ has characteristic $0$ then $S$ is a DVR, and if $A$ has
characteristic $p^\nu$ then $S$ is an Artinian ring.

\begin{theorem}\label{thm:29.6}
Let $(A, m, K)$ be a local ring, and suppose $\chark K = p$. Let $C \subset A$ be a
subring, and assume that $C$ is a Noetherian local ring with maximal ideal $pC$,
and that $K = A/m$ is separable over $C/pC$. Then there exists a quasi-coefficient
ring $S$ of $A$ containing $C$; moreover, if $A$ is flat over $C$, then it is also
flat over $S$.
\end{theorem}

\begin{proof}
Let $\{B_\lambda\}_{\lambda \in \Lambda}$ be a $p$-basis of $K$ over $C/pC$, and choose an inverse
image $b_\lambda \in A$ for each $\beta_\lambda$. Setting $C[\{b_\lambda\}] = C'$, by
Theorem~\hyperref[thm:26.8]{26.8} we see that $C'/(m \cap C') = (C/pC)[\{\bar{X}_\lambda\}]$ is a
polynomial ring over $C/pC$. Hence if $f(\ldots, X_\lambda, \ldots)$ is a non-zero
polynomial with coefficients in $C$ which satisfies $f(b_\lambda) \in m$, then setting
$\rho$ for the highest common factor of the coefficients of $f$, we have
\[
  f(X) = \rho \cdot f_1(X) \quad \text{and} \quad f_1(b_\lambda) \not\equiv 0.
\]
Thus $f_1(b_\lambda) \notin m$, and we must have $\rho \in m$, so that we have shown that
$m \cap C' = pC'$. Setting $S = (C')_{m}$, we have $S \subset A$, $m \cap S = pS$ and
$S/pS = K'$; since $C'$ is $p$-adically separated, so is $S$, and hence all the
ideals of $S$ are of the form $(0)$ or $(p^\nu)$. Thus $S$ is Noetherian, and it
satisfies all the conditions for a quasi-coefficient ring of $A$. If $A$ is flat over
$C$, then
\[
  p^\nu C \otimes_C A \cong p^\nu A.
\]
and hence the composite
\[
  p^\nu C \otimes A = (p^\nu C \otimes_C S) \otimes_S A \to p^\nu S \otimes_S A \to p^\nu A
\]
is injective; but the first arrow is surjective, so that the second arrow
$p^\nu S \otimes_S A \to p^\nu A$ is injective. By Theorem~\hyperref[thm:7.7]{7.7}, this proves that $A$ is
flat over $S$.
\end{proof}

All the assumptions of this theorem are satisfied by taking $C$ to be the image
of $\mathbb{Z}_{(p)} \to A$, so that this proves that every local ring has a quasi-coefficient
ring (including the quasi-coefficient field of a local ring in the equal
characteristic case).

\begin{theorem}\label{thm:29.7}
Let $(A, m, K)$ be a local ring, and $\hat{A}$ its completion, and suppose
$\chark K = p$. Let $S$ be a quasi-coefficient ring of $A$, and write $\bar{S}$ for its
image in $\hat{A}$; then there exists a unique coefficient ring $A_0$ of $\hat{A}$
containing $\bar{S}$.
\end{theorem}

\begin{proof}
Since $\bar{S}$ is a quasi-coefficient ring of $\hat{A}$, we can assume that $A$ is a
complete local ring. If $A$ has characteristic $0$ then $S$ is a DVR, and by
Theorem~\hyperref[thm:29.1]{1} there exists a complete $p$-ring $R$ containing $S$ and with
residue field $K$. Now $K$ is 0-\'etale over $S/pS$, and $R$ is flat over $S$, so that
$R$ is $pR$-\'etale over $S$ (by Ex.~\ref{ex:28.1-ex}). Hence there exists a unique
$S$-algebra homomorphism $R \to A$ which induces the identity map on the residue
fields; we write $A_0$ for the image of $R$. Then $R \cong A_0$, and $A_0$ is a
coefficient ring of $A$. If $A$ had another coefficient ring $B$ containing $S$ then
$B$ would also be a complete $p$-ring, and for the same reason as above, there
exist unique $S$-algebra homomorphisms $B \to A_0$ and $B \to A_0$, so that we must
have $B = A_0$.

If $A$ has characteristic $p^\nu$ then $S$ is an Artinian ring, and therefore
complete, so that applying Theorem~\hyperref[thm:29.4]{4} to $S$ itself, we can write
$S = R_0/p^\nu R_0$ with $R_0$ a complete $p$-ring. Let $R$ be a complete $p$-ring
containing $R_0$ and with residue field $K$; then by Ex.~\ref{ex:28.1-ex}, $R$ is $pR$-\'etale
over $R_0$, so that there exists a unique $R_0$-algebra homomorphism $R \to A$
inducing the identity map on the residue field $K$. The image is a coefficient
ring of $A$ containing $S$; the uniqueness is proved as in the case of
characteristic $0$.
\end{proof}

Next we study the structure of complete regular local rings. Let $(A, m, K)$ be a
local ring of unequal characteristic, and suppose that $\chark K = p$; then $A$ is
said to be \emph{ramified}\index{ramified} if $p \in m^2$ and \emph{unramified} if $p \notin m^2$. We will
also say that $A$ is unramified in the case of equal characteristic.

\begin{theorem}\label{thm:29.8}
An unramified complete regular local ring is a formal power series ring over a
field or over a complete $p$-ring.
\end{theorem}

\begin{proof}
Let $R$ be a coefficient ring of $A$. In the case of equal characteristic, $R$ is a
field, and if $x_1, \ldots, x_n$ is a regular system of parameters of $A$ then
$A = R[x_1, \ldots, x_n] \cong R[[X_1, \ldots, X_n]]$ (see the proof of
Theorem~\hyperref[thm:29.5]{5}). In the case of unequal characteristic, $R$ is a complete $p$-ring,
and since $p \notin m^2$, we can choose a regular system of parameters
$\{p, x_2, \ldots, x_n\}$ of $A$ containing $p$. Then $A = R[x_2, \ldots, x_n] \cong
R[[X_2, \ldots, X_n]]$.
\end{proof}

In the ramified case, it is not necessarily the case that $A$ can be expressed as
a formal power series ring over a DVR. To give the structure theorem in this
case we need the notion of an \emph{Eisenstein extension}.

\begin{lemma}[\textbf{Eisenstein's irreducibility criterion}]\label{lem:29.1}
Let $A$ be a ring, and $f(X) = X^n + a_1X^{n-1} + \cdots + a_n$, with $a_i \in A$. If there
exists a prime ideal $P$ of $A$ such that $a_1, \ldots, a_n \in P$ but $a_n \notin P^2$ then $f$
is irreducible in $A[X]$. If in addition $A$ is an integrally closed domain then
the principal ideal $(f)$ is a prime ideal of $A[X]$.
\end{lemma}

\begin{proof}
If $f$ is reducible, we can write
\[
  f = (X^r + b_1X^{r-1} + \cdots + b_r)(X^s + c_1X^{s-1} + \cdots + c_s)
\]
with $0 < r < n$, $s = n - r$ and $b_i, c_j \in A$. Reducing the coefficients on either
side modulo $P$ we have
\[
  X^n = (X^r + \bar{b}_1X^{r-1} + \cdots + \bar{b}_r)(X^s + \bar{c}_1X^{s-1} + \cdots + \bar{c}_s)
\]
in $(A/P)[X]$, so that we must have $b_i, c_j \in P$ for all $i, j$, but then
$a_n = b_rc_s \in P^2$, which contradicts the assumption. If $A$ is an integrally closed
domain and $K$ is the field of fractions of $A$, then by Ex.~\ref{ex:9.6-ex}, $f$ remains
irreducible in $K[X]$. Also, $f$ is manic, so that we have
$f \cdot A[X] = f \cdot K[X] \cap A[X]$, and this is a prime ideal of $A[X]$.
\end{proof}

Let $(A, m)$ be a normal local ring; then an extension ring
\[
  B = A[x]/(f) = A[x]
\]
defined by an \emph{Eisenstein polynomial}\index{Eisenstein polynomial}
\[
  f = X^n + a_1X^{n-1} + \cdots + a_n \quad \text{with } a_i \in m \text{ for all } i,
      \text{ and } a_n \notin m^2
\]
is called an \emph{Eisenstein extension} of $A$. By the lemma, $B$ is an integral domain,
and is integral over $A$. We have $B/mB = (A/m)[X]/(X^n)$, so that $B$ has just one
maximal ideal $n = mB + xB$. Hence $B$ is a local ring, and its residue field
coincides with that of $A$.

\begin{theorem}\label{thm:29.9}
\begin{enumerate}[(i)]
  \item If $(A, m)$ is a regular local ring, then an Eisenstein extension of $A$ is
        again a regular local ring.
  \item If $A$ is a ramified complete regular local ring and $R$ is a coefficient
        ring of $A$ then there is a subring $A_1 \subset A$ with the following
        properties:
        \begin{enumerate}[(1)]
          \item $A_1$ is an unramified complete regular local ring containing $R$,
                and hence can be expressed as a formal power series ring over $R$;
          \item $A$ is an Eisenstein extension of $A_1$.
        \end{enumerate}
\end{enumerate}
\end{theorem}

\begin{proof}
(i) Let $B = A[x]$, and $x^n + a_1x^{n-1} + \cdots + a_n = 0$, with $a_i \in m$ and
$a_n \in m^2$ but $a_n \notin m^3$. Then there exists a regular system of parameters
$\{y_1, \ldots, y_n = a_n\}$ of $A$ with $a_n$ as an element. As we have seen above, the
maximal ideal of $B$ is $mB + xB$, but $a_n \in xB$, so that
$\{y_1, \ldots, y_{n-1}, x\}$ is a regular system of parameters of $B$.

(ii) Since $\htop pA = 1$, by a skilful choice of a regular system of parameters
$\{x_1, \ldots, x_d\}$ of $A$, we can arrange that $\{p, x_2, \ldots, x_d\}$ is a system of
parameters of $A$. If we set $A_1 = R[x_2, \ldots, x_d]$ then $A_1$ is a complete
unramified regular local ring, and $A$ is a finite module over $A_1$ (see the
proof of Theorem~\hyperref[thm:29.5]{5}). We set $m_1$ for the maximal ideal of $A_1$. Now
$A = m_1A + A_1[x_1]$, so that by Theorem~\hyperref[thm:8.4]{8.4} (or by NAK),
$A = A_1[x_1]$. Let
\[
  f(X) = X^n + a_1X^{n-1} + \cdots + a_n \quad \text{with } a_i \in A_1
\]
be the minimal polynomial of $x_1$ over $A_1$. Then $a_n = x_1 \cdot \ldots \in m$,
so that $a_n \notin m_1^2$. Therefore by Hensel's lemma (Theorem~\hyperref[thm:8.3]{8.3}), all the
$a_i \in m_1$. We are left to prove that $a_n \notin m_1^2$. Write $p = \sum_{i=1}^{d} b_ix_i$ with
$b_i \in A_1$, and express the $b_i$ in the form $b_i = q_i(x_1)$, with
$q_i(X) \in A_1[X]$; then $x_1$ is a root of
\[
  F(X) = q_1(X)X + \sum_{i=2}^{d} q_i(X)x_i - p
\]
so that $F(X)$ is divisible by $f(X)$. Hence the constant term $F(0)$ of $F$ is
divisible by $a_n$. However, $F(0) = \sum_{i=2}^{d} q_i(0)x_i - p$, and
$p, x_2, \ldots, x_d$ is a regular system of parameters, so that $F(0) \notin m^2$, hence
also $a_n \notin m_1^2$.
\end{proof}

% ============ 习题 ============
\section*{Exercises to \S\ref{sec:29}}
\emph{Prove the following propositions.}

\begin{exercise}\label{ex:29.1-ex}
Let $A$ be a complete $p$-ring, $y$ an indeterminate over $A$, and $B = A[[y]]$; let
$C = B[x]$ be the Eisenstein extension of $B$ given by $x^2 + yx + p = 0$. Then $C$
is a two-dimensional complete regular local ring, but is not a formal power
series ring over a DVR of characteristic $0$.
\end{exercise}

\begin{exercise}\label{ex:29.2-ex}
In Theorem~\hyperref[thm:29.2]{29.2}, if $k$ is a perfect field then $\varphi$ is uniquely determined
by $\varphi_0$.
\end{exercise}

% ============ §30 Connections with derivations (原书 pp.230-245) ============
\section{Connections with derivations}\label{sec:30}\index{Jacobian criterion, section 30}\index{weak Jacobian condition, section 30}\index{Nagata-Zariski-Lipman, section 30}\index{geometrically regular, section 30}

\begin{theorem}[\textbf{Nagata-Zariski-Lipman}]\label{thm:30.1}
Let $(A, m)$ be a complete Noetherian local ring with $\mathbb{Q} \subset A$. Suppose that
$x_1, \ldots, x_r \in m$ and $D_1, \ldots, D_r \in \Der(A)$ are elements satisfying
$\det(D_ix_j) \notin m$. Then
\begin{enumerate}[(i)]
  \item There is a subring $C \subset A$ such that
        \[
          A = C[x_1, \ldots, x_r] \cong C[[X_1, \ldots, X_r]].
        \]
        Therefore $x_1, \ldots, x_r$ are analytically independent over $C$, and $A$ is
        $I$-smooth over $C$, where $I = \sum Ax_i$, and therefore also $m$-smooth over
        $C$.
  \item If $\mathfrak{g} = \sum AD_i$ is a Lie algebra, (that is if $[D_i, D_j] \in \mathfrak{g}$ for all
        $i, j$) then we can take $C$ to be $\{a \in A \mid D_1a = \cdots = D_ra = 0\}$.
\end{enumerate}
\end{theorem}

\begin{proof}
Letting $(c_{ij})$ be the inverse matrix of $(D_ix_j)$, and setting
$D_i' = \sum c_{ij}D_j$, we have $D_i'x_j = \delta_{ij}$, so that we can assume that
$D_ix_j = \delta_{ij}$. Quite generally, for an element $t \in m$ and a derivation
$D \in \Der(A)$, we define a map $E(D, t)\colon A \to A$ by
\[
  E(D, t) = \sum_{n \ge 0} \frac{t^n}{n!}D^n;
\]
by our assumptions, $E(D, t)(a) = \sum (t^n/n!)D^n(a)$ is meaningful, and one sees
easily that $E(D, t)$ is a ring homomorphism. Now set
\[
  E_i = E(D_i, -x_i) \quad \text{and} \quad C_i = \Im(E_i);
\]
then $C_i$ is a subring of $A$, and by computation we see that
\[
  D_i\Big(\sum \frac{(-x_i)^n}{n!}D_i^n(a)\Big) = \sum \frac{(-x_i)^{n-1}}{(n-1)!}D_i^n(a)
      + \sum \frac{(-x_i)^n}{n!}D_i^{n+1}(a) = 0,
\]
so that $C_i \subset \{a \in A \mid D_ia = 0\}$. Conversely, if $D_ia = 0$ then
$E_i(a) = a$, so that $C_i = \{a \in A \mid D_ia = 0\}$. Also, since for any $a \in A$ we have
$E_i(a) \equiv a \mod x_1A$, we see that elements of $A$ can be expanded in power
series in $x_1$ with coefficients in $C_1$, so that $A = C_1[x_1]$. Now
$E_1(x_1) = x_1 - x_1 = 0$, so that $x_1A \subset \Ker E_1$, and conversely, if
$E_1(a) = a - x_1D_1a + \cdots = 0$ then $a \in x_1A$, and therefore $\Ker E_1 = x_1A$.
Also, $D_1 \circ D_1c = 0 \iff E_1(c) = c$, so that $C_1 \cap x_1A = 0$. Now we prove
that $x_1$ is analytically independent over $C_1$; by contradiction, suppose that
\[
  c_\nu x_1^\nu + c_{\nu+1}x_1^{\nu+1} + \cdots = 0 \quad \text{with } c_\nu \in C_1
      \text{ and } c_\nu \ne 0.
\]
Then by Ex.~\hyperref[ex:25.2-ex]{25.2}, since $x_1$ is not a zero-divisor in $A$, we have
$c_\nu \in x_1A$, which is a contradiction. Thus if $0 \ne \varphi(X) \in C_1[X]$ then
$\varphi(x_1) \ne 0$, as required.

If $r > 1$, then write $D_i'$ for the restriction to $C_1$ of $E_1 \circ D_i$; then
$D_i' \in \Der(C_1)$, and $x_j \in C_1$ with $D_i'x_j = \delta_{ij}$ for $2 \le i, j \le r$, so that by
induction we have
\[
  C_1 = C[x_2, \ldots, x_r] \cong C[[X_2, \ldots, X_r]].
\]
(i) follows from this.

If $\mathfrak{g}$ is a Lie algebra, then we first arrange as before that $D_ix_j = \delta_{ij}$,
and then set $[D_i, D_j] = \sum_\nu a_{ij\nu}D_\nu$ with $a_{ij\nu} \in A$; then
$[D_i, D_j]x_\nu = D_i(\delta_{j\nu}) - D_j(\delta_{i\nu}) = 0$, so that $a_{ij\nu} = 0$, hence
$[D_i, D_j] = 0$, and $D_1(D_i(C_1)) = D_i(D_1(C_1)) = D_i(0) = 0$. Therefore
$D_i(C_1) \subset C_1$ for $i > 1$, and then in the above notation $D_i' = D_i$ for $i > 1$.
Thus by induction we have $C = \{a \in A \mid D_1a = \cdots = D_ra = 0\}$.
\end{proof}

\begin{corollary*}
Let $(A, m)$ be a reduced $n$-dimensional local ring containing $\mathbb{Q}$, and suppose
that the completion $\hat{A}$ of $A$ is also reduced. If there exist $n$ elements
$D_1, \ldots, D_n \in \Der(A)$ and $x_1, \ldots, x_n \in m$ such that
$\det(D_ix_j) \notin m$ then $A$ is a regular local ring and $x_1, \ldots, x_n$ is a regular
system of parameters of $A$. Suppose in addition that $\mathfrak{g} = \sum AD_i$ is a Lie
algebra; then $C = \{a \in A \mid D_1a = \cdots = D_na = 0\}$ is a coefficient field of
$\hat{A}$.
\end{corollary*}

\begin{proof}
Consider $\hat{A}$, with each of the $D_i$ extended to $\hat{A}$. If
$[D_i, D_j] = \sum_\nu a_{ij\nu}D_\nu$ holds in $A$ then it also holds in $\hat{A}$, so that if
$\mathfrak{g}$ is a Lie algebra, so is $\hat{\mathfrak{g}} = \sum \hat{A}D_i$. By the theorem,
$\hat{A} = C[[x_1, \ldots, x_n]]$, and $C$ is isomorphic to $\hat{A}/\sum x_i\hat{A}$, so is a
zero-dimensional local ring. Now by assumption $C$ is also reduced, so that $C$
must be a field. Therefore $\hat{A}$ is a formal power series ring over a field, and
hence is regular, so that $A$ is also regular. The other assertion is also clear.
\end{proof}

\begin{remark}
If we view this corollary as a criterion for regularity, then the condition that
$\hat{A}$ should be reduced is rather a nuisance; however, as we will see later, for
a very wide variety of local rings, we have $A$ is reduced $\iff \hat{A}$ is reduced.
This is the case (corollary of Theorem~\hyperref[thm:32.6]{32.6}) if $A$ is a localisation of a ring
$B$ which is finitely generated over a field (such an $A$ is said to be
\emph{essentially of finite type} over $K$). Note also that if we start off with a
regular local ring $A$, then the corollary gives a concrete method of
constructing a coefficient field of $\hat{A}$.
\end{remark}

Next we consider rings which are finitely generated over a field, which are
important in algebraic geometry.

\begin{theorem}\label{thm:30.2}
Let $k$ be a field, and $A = k[x_1, \ldots, x_n]$ a finitely generated ring over $k$. If
$A_p$ is 0-smooth over $k$ for every $p \in \mathrm{m}$-$\Spec A$, then $A$ is 0-smooth over
$k$.
\end{theorem}

\begin{proof}
Write $k[X]$ for $k[X_1, \ldots, X_n]$, and let $I = \{f(X) \in k[X] \mid f(x) = 0\}$, so
that $A = k[X]/I$. Suppose that $I = (f_1, \ldots, f_s)$. Consider a commutative diagram
\[
\begin{tikzcd}
A \arrow[r, "\bar{\varphi}"] & C/N \\
k \arrow[u] \arrow[r] & C \arrow[u]
\end{tikzcd}
\]
where $C$ is a ring, and $N \subset C$ is an ideal satisfying $N^2 = 0$. To lift
$\bar{\varphi}$ to $A \to C$, we first of all choose $u_i \in C$ such that
$\bar{\varphi}(x_i) = \bar{u}_i$. If $f \in I$ then $f(u_1, \ldots, u_n) \in N$. Now if we can choose
$y_i \in N$ for $1 \le i \le n$ such that $f_j(u + y) = 0$ for all $j$, then the
homomorphism $A \to C$ defined by $X_i \mapsto u_i + y_i$ is a lifting of
$\bar{\varphi}$. We have
\[
  f_j(u + y) = f_j(u) + \sum_{i=1}^{n} (\partial f_j/\partial X_i)(u) \cdot y_i,
\]
so that we are looking for solutions in $N$ of the system of linear equations in
$y_1, \ldots, y_n$:
\[
  (*) \quad f_j(u) + \sum_{i=1}^{n} (\partial f_j/\partial X_i)(u) \cdot y_i = 0
      \quad \text{for } j = 1, \ldots, s.
\]
For each maximal ideal $p$, the local ring $A_p$ is 0-smooth over $k$, so that if we
set $S = \bar{\varphi}(A - p)$ and $\tilde{S} = \varphi^{-1}(S)$ then in the diagram
\[
\begin{tikzcd}
A_p \arrow[r, "\bar{\varphi}_p"] & C_{\tilde{S}}/N_{\tilde{S}} = C_S/N_S \\
k \arrow[u] \arrow[r] & C_{\tilde{S}} \arrow[u]
\end{tikzcd}
\]
there exists a $y^p\colon A_p \to C_{\tilde{S}}$ lifting $\bar{\varphi}_p$. From this we see that (*), as
a system of equations in $N_{\tilde{S}}$, has a solution in $N_{\tilde{S}}$. If we view $N$ as a
$C/N$-module then $N_{\tilde{S}} = N_S$. Thus the theorem reduces to the following
lemma.
\end{proof}

\begin{lemma}\label{lem:30.1}
Let $A$ and $B$ be rings, $\varphi\colon A \to B$ a ring homomorphism, and $N$ a
$B$-module; suppose that $b_{ij} \in B$ and $\beta_i \in N$. If the system of linear
equations
\[
  \sum_j b_{ij}y_j = \beta_i \quad (\text{for } i = 1, \ldots, s)
\]
has a solution in $N_{A-p}$ for every $p \in \mathrm{m}$-$\Spec A$, then it has a solution in $N$.
\end{lemma}

\begin{proof}
The assumption that there is a solution in $N_{A-p}$ means that there exist
$q_{jp} \in N$ for $1 \le j \le n$, and $t_p \in A - p$ such that
\[
  \sum_j b_{ij}q_{jp} - \varphi(t_p)\beta_i = 0 \quad \text{for } 1 \le i \le s.
\]
Now since $\sum_p t_pA = A$, there is a finite set $\{p_1, \ldots, p_\ell\} \subset
\mathrm{m}$-$\Spec A$ and $a_\nu \in A$ such that
\[
  \sum_{\nu=1}^{\ell} a_\nu t_{p_\nu} = 1.
\]
Hence if we set
\[
  q_j = \sum_\nu \varphi(a_\nu)q_{jp_\nu},
\]
we get $\sum_j b_{ij}q_j = \beta_i$ for $1 \le i \le s$.
\end{proof}

\begin{theorem}\label{thm:30.3}
Let $k$ be a field, $S = k[X_1, \ldots, X_n]$, and $I$, $P$ ideals of $S$ such that
$I \subset P \in \Spec S$. Set
\[
  S_P = R, \quad \rad(R) = PR = M, \quad \text{and} \quad R/IR = A, \quad \rad(A) = m,
\]
\[
  R/M = A/m = K,
\]
and suppose that $\htop IR = r$ and $I = (f_1(X), \ldots, f_t(X))$. Then the following
conditions are equivalent.
\begin{enumerate}[(1)]
  \item $\rank(\partial f_i/\partial X_j \bmod P)_{i=1, \ldots, t;\, j=1, \ldots, n} = r$;
  \item $A$ is 0-smooth over $k$;
  \item $A$ is $m$-smooth over $k$;
  \item $\Omega_{A/k}$ is a free $A$-module of rank $n - r$;
  \item $A$ is an integral domain, its field of fractions is separable over $k$, and
        $\Omega_{A/k}$ is a free $A$-module.
\end{enumerate}
\end{theorem}

\begin{proof}
Note that $R$ is a regular local ring.

(1) $\Rightarrow$ (2) By assumption, for a suitable choice of $r$ elements
$D_1, \ldots, D_r$ from $\partial/\partial X_1, \ldots, \partial/\partial X_n$, and of $r$ elements
$g_1, \ldots, g_r$ from $f_1, \ldots, f_t$, we have $\det(D_ig_j) \notin M$. We observe that
taking $R \ni f$ into $(\partial f/\partial X_1 \bmod M, \ldots, \partial f/\partial X_n \bmod M)
\in K^n$ induces a linear map $M/M^2 \to K^r$, so that the images of $g_1, \ldots, g_r$ in
$M/M^2$ are linearly independent. Therefore $\sum g_iR$ is a height $r$ prime ideal
contained in $IR$, and hence $\sum g_iR = IR$. Given a commutative diagram
\[
\begin{tikzcd}
A \arrow[r, "\bar{\varphi}"] & C/N \\
k \arrow[u] \arrow[r] & C \arrow[u]
\end{tikzcd}
\]
with $N^2 = 0$, write $\bar{X}_i \in A$ for the image of $X_i$, and choose $u_i \in C$
such that $\varphi(u_i) = \bar{\varphi}(\bar{X}_i) \in C/N$. Then there is a homomorphism
$R \to C$ defined by $X_i \mapsto u_i$, and this induces a lifting of $\bar{\varphi}$ to $C$ if and
only if $g_i(u) = 0$ for $1 \le i \le r$. Therefore, as in the proof of the previous
theorem, we need only solve the system of equations in unknowns
$y_1, \ldots, y_n \in N$:
\[
  (*) \quad g_i(u) + \sum_{j=1}^{n} (\partial g_i/\partial X_j)(u) \cdot y_j = 0
      \quad \text{for } 1 \le i \le r.
\]
However, we view $N$ as a $C/N$-module, then via $\bar{\varphi}$ as an $A$-module, so that we
can replace the coefficients $(\partial g_i/\partial X_j)(u)$ of this system by
$(\partial g_i/\partial X_j)(\bar{X}) \in A$. Now by assumption, one $r \times r$ minor of this
$r \times n$ matrix of coefficients is a unit of $A$, so that (*) can always be solved.

(2) $\Rightarrow$ (3) is trivial.

(3) $\Rightarrow$ (1) By \S\ref{sec:28}, Lemma~\hyperref[lem:28.1]{1}, $A$ is regular, so that $IR$ is generated by
$r$ elements forming a subset of a regular system of parameters of $R$, and the
image of the natural map $IR \to M \to M/M^2$ is an $r$-dimensional $K$-vector space.
Let $k_1 \subset k$ be the prime subfield; then by
Theorem~\hyperref[thm:26.9]{26.9} $K = R/M$ is 0-smooth over $k_1$, so that by
Theorem~\hyperref[thm:25.2]{25.2}, the sequence
\[
  0 \to M/M^2 \to \Omega_{R/k_1} \otimes_R K \to \Omega_{K/k_1} \to 0
\]
is exact. We can write $\Omega_{S/k_1} = (\Omega_{k/k_1} \otimes S) \oplus F$, where $F$ is the free
$S$-module with basis $dX_1, \ldots, dX_n$ (for example, by
Theorem~\hyperref[thm:25.1]{25.1}), and localising we get
\[
  \Omega_{R/k_1} = (\Omega_{k/k_1} \otimes R) \oplus F \otimes_S R.
\]
Hence $\Omega_{R/k_1} \otimes_R K = (\Omega_{k/k_1} \otimes_k K) \oplus (F \otimes K)$. However, from
\[
  0 \to I/I^2 \to \Omega_{S/k_1} \otimes_S A \to \Omega_{A/k_1} \to 0
\]
we get the exact sequence $(I/I^2) \otimes_A K \to \Omega_{R/k_1} \otimes_R K \to
\Omega_{A/k_1} \otimes_A K \to 0$, and if $A$ is $m$-smooth over $k$ then by the corollary to
Theorem~\hyperref[thm:28.6]{28.6}, $\Omega_{k/k_1} \otimes K \to \Omega_{A/k_1} \otimes K$ is injective, so that
$V = \Im((I/I^2) \otimes K \to \Omega_{R/k_1} \otimes K)$ maps isomorphically to its
projection $W \subset F \otimes K$ in the second factor of the decomposition
$\Omega_{R/k_1} \otimes K = (\Omega_{k/k_1} \otimes K) \oplus (F \otimes K)$. Now factor
$I/I^2 \otimes K \to \Omega_{R/k_1} \otimes K$ as the composite
$I/I^2 \otimes K \to M/M^2 \to \Omega_{R/k_1} \otimes K$; as we have seen above, the first arrow
has rank $r$, and the second is injective, so that $\rank V = r$. Now
$F \otimes K = KdX_1 + \cdots + KdX_n$, and if we write
$(\partial f_i/\partial X_j) \bmod P = a_{ij}$ then $W$ is spanned by $\sum_j a_{ij}dX_j$ for
$1 \le i \le t$. Therefore $\rank(a_{ij}) = r$, and this proves (1).

(2) $\Rightarrow$ (4) By Theorem~\hyperref[thm:25.2]{25.2},
\[
  0 \to IR/I^2R \to \Omega_{R/k} \otimes_R A \to \Omega_{A/k} \to 0
\]
is a split exact sequence, and since $\Omega_{R/k} \otimes_R A$ is a free $A$-module with basis
$dX_1, \ldots, dX_n$, the $A$-module $\Omega_{A/k}$ is projective; but $A$ is a local ring, so
that $\Omega_{A/k}$ is free. Also $A$ is a regular local ring, therefore an integral
domain, and if $L$ is its field of fractions then $L$ is 0-smooth over $A$, hence also
0-smooth over $k$, so that $L$ is separable over $k$.

(4) $\Rightarrow$ (5) By Theorem~\hyperref[thm:26.2]{26.2}, the field of fractions $L$ of $A$ is
separably generated over $k$, so that a separating transcendence basis of $L$ over
$k$ is a differential basis, and $\rank \Omega_{L/k} = \trdeg_k L = n - r$; but
$\Omega_{L/k} = \Omega_{A/k} \otimes_A L$ and hence $\rank \Omega_{A/k} = \rank \Omega_{L/k} = n - r$.

(5) $\Rightarrow$ (1) In the exact sequence
\[
  IR/I^2R \to \Omega_{R/k} \otimes_R A \to \Omega_{A/k} \to 0,
\]
set $E = \Im(IR/I^2R \to \Omega_{R/k} \otimes_R A)$. Then
$\Omega_{R/k} \otimes_R A \cong E \oplus A^{n-r}$, so that $E \cong A^r$, and therefore
$E \otimes K \cong K^r$. This gives (1).
\end{proof}

\begin{remark}
In the above proof, the equivalence of (1), (2), (4) and (5) was comparatively
easy. The proof of (3) $\Rightarrow$ (1) used the corollary to Theorem~\hyperref[thm:28.6]{28.6}, and so
is not very elementary.
\end{remark}

\begin{remark}
If $k$ is a perfect field, or more generally if the residue field $K$ is separable
over $k$, then $m$-smoothness is equivalent to $A$ being a regular local ring, so
that Theorem~\hyperref[thm:30.3]{3} gives a criterion for regularity. In the case of an imperfect
field $k$, if $A$ is 0-smooth over $k$ then so is the field of fractions $L$ of $A$, but
the residue field $K$ is not necessarily separable over $k$; for example,
$A = k[X]_{(X^p - a)}$ with $a \in k - k^p$. For the case of an imperfect field $k$, we
give a regularity criterion for $A$ in Theorem~\hyperref[thm:30.5]{5}.
\end{remark}

Quite generally, let $A$ be a ring and $P$ a prime ideal of $A$, and let
$D_1, \ldots, D_s \in \Der(A)$ and $f_1, \ldots, f_t \in A$; then we write
\[
  J(f_1, \ldots, f_t; D_1, \ldots, D_s)(P) = (D_if_j \bmod P).
\]
This is an $s \times t$ matrix with entries in the integral domain $A/P$.

\begin{theorem}\label{thm:30.4}
Let $R$ be a regular ring, $P \in \Spec R$, and let $I \subset P$ be an ideal of $R$;
suppose that $\htop IR_P = r$.
\begin{enumerate}[(i)]
  \item for any $D_1, \ldots, D_s \in \Der(R)$ and $f_1, \ldots, f_t \in I$ we have
        \[
          \rank J(f_1, \ldots, f_t; D_1, \ldots, D_s)(P) \le r;
        \]
  \item if $D_1, \ldots, D_r \in \Der(R)$ and $f_1, \ldots, f_r \in I$ are such that
        $\det(D_if_j) \notin P$ then $IR_P = (f_1, \ldots, f_r)R_P$ and $R_P/IR_P$ is regular.
\end{enumerate}
\end{theorem}

\begin{proof}
(i) If $Q$ is a prime ideal of $R$ with $I \subset Q \subset P$ and $\htop Q = r$ then
\[
  \rank J(f_1, \ldots; D_1, \ldots)(P) \le \rank J(f_1, \ldots; D_1, \ldots)(Q),
\]
and if we set $QR_Q = \mathfrak{n}$ then $R_Q$ is an $r$-dimensional regular local ring, so
that $m$ is generated by $r$ elements, $m = (g_1, \ldots, g_r)$. Working in $R_Q$, we can
write
\[
  f_j = \sum_{\nu=1}^{r} g_\nu a_{\nu j} \quad \text{with } a_{\nu j} \in R_Q,
\]
so that
\[
  D_if_j \equiv \sum_{\nu=1}^{r} (D_ig_\nu) \cdot a_{\nu j} \bmod Q,
\]
and therefore
\[
  \rank J(f_1, \ldots, f_t; D_1, \ldots, D_s)(P) \le
  \rank J(g_1, \ldots, g_r; D_1, \ldots, D_s)(Q) \le r.
\]
(ii) Set $M = PR_P$; then if $\det(D_if_j) \notin M$ one sees easily that the images in
$M/M^2$ of $f_1, \ldots, f_r$ are linearly independent over $R_P/M = \kappa(M)$, so that
$\sum f_iR_P$ is a prime ideal of height $r$, and therefore coincides with $IR_P$.
Also, $R_P/IR_P$ is regular.
\end{proof}

\begin{theorem}[\textbf{Zariski}]\label{thm:30.5}
Let $k$ be a field of characteristic $p$, $S = k[X_1, \ldots, X_n]$, and $I$ and $P$
ideals of $S$ with $I \subset P \in \Spec S$. Set $S_P = R$ and $R/IR = A$, and suppose
that $\htop IR = r$ and $I = (f_1, \ldots, f_t)$. Then the following three conditions are
equivalent.
\begin{enumerate}[(1)]
  \item $A$ is a regular local ring.
  \item For any $p$-basis $\{u_\gamma\}_{\gamma \in \Gamma}$ of $k$, define $D_\gamma \in \Der(S)$ by
        $D_\gamma(u_{\gamma'}) = \delta_{\gamma\gamma'}$ (the Kronecker $\delta$) and
        $D_\gamma X_i = 0$; then there are a finite number of elements
        $\alpha, \beta, \ldots, \gamma \in \Gamma$ such that
        \[
          \rank J(f_1, \ldots, f_t; D_\alpha, D_\beta, \ldots, D_\gamma,
              \partial/\partial X_1, \ldots, \partial/\partial X_n)(P) = r.
        \]
  \item There exists a subfield $k' \subset k$ with the following properties:
        $k^p \subset k' \subset k$, $[k : k'] < \infty$, and $\Omega_{A/k'}$ is a free $A$-module, with
        \[
          \rank \Omega_{A/k'} = n - r + \rank \Omega_{k/k'}.
        \]
\end{enumerate}
\end{theorem}

\begin{proof}
(2) $\Rightarrow$ (1) comes from Theorem~\hyperref[thm:30.4]{4}, (ii).

(1) $\Rightarrow$ (2), (3) If $A$ is regular then $IR$ is generated by $r$ elements, and
these form an $R$-sequence, so that $IR/I^2R$ is a free $A$-module of rank $r$. Set
$M = PR$, $m = M/IR$ and $K = R/M = A/m$; then the image of $IR \to M/M^2$ is an
$r$-dimensional $K$-vector space, so that the natural map $(IR/I^2R) \otimes_A K \to
M/M^2$ is injective. From the exact sequence
\[
  IR/I^2R \to \Omega_{R/k} \otimes_R A \to \Omega_{A/k} \to 0
\]
we get the exact sequence
\[
  (*) \quad (IR/I^2R) \otimes_A K \to \Omega_{R/k} \otimes_R K \to \Omega_{A/k} \otimes_A K \to 0.
\]
Now by Theorem~\hyperref[thm:25.2]{25.2},
\[
  0 \to M/M^2 \to \Omega_{R/k} \otimes_R K \to \Omega_{K/k} \to 0
\]
is an exact sequence. The first arrow in (*) is the composite
$(IR/I^2R) \to M/M^2 \to \Omega_{R/k} \otimes K$, and so is injective. Thus
\[
  (**) \quad 0 \to (IR/I^2R) \otimes K \to \Omega_{R/k} \otimes K \to \Omega_{A/k} \otimes K \to 0
\]
is exact. Let $\{du_\sigma\}_{\sigma \in \Sigma}$ be a basis of $\Omega_k$ over $k$; then $\Omega_S$ is the free
$S$-module with basis $\{du_\sigma\}_{\sigma \in \Sigma} \cup \{dX_1, \ldots, dX_n\}$, and
$\Omega_R = \Omega_S \otimes_S R$, $\Omega_R \otimes_R K = \Omega_k \otimes_k K \oplus
(KdX_1 \oplus \cdots \oplus KdX_n)$. Now reorder $f_1, \ldots, f_t$ so that $f_1, \ldots, f_r$ are
generators of $IR$; then $df_1, \ldots, df_r \in \Omega_R \otimes K$ can be expressed using
$dX_1, \ldots, dX_n$ together with finitely many elements $du_{\tau_1}, \ldots, du_{\tau_s}$, and this
gives (2). Now if we let $k'$ be the field obtained by adjoining to $k^p$ all the
elements $u_\sigma$ for $\sigma \in \Sigma$ other than the $du_{\tau_1}, \ldots, du_{\tau_s}$ just used, then
from (**) we see that
\[
  0 \to (IR/I^2R) \otimes K \to \Omega_{R/k'} \otimes K \to \Omega_{A/k'} \otimes K \to 0
\]
is also exact. In the exact sequence
\[
  (\dagger) \quad IR/I^2R \to \Omega_{R/k'} \otimes_R A \to \Omega_{A/k'} \to 0,
\]
the middle term is the free $A$-module with basis
$du_{\tau_1}, \ldots, du_{\tau_s}, dX_1, \ldots, dX_n$; now the generators $f_1, \ldots, f_r$ of $IR$ map
to $df_1, \ldots, df_r \in \Omega_{R/k'} \otimes_R A$, and since these are linearly independent over
$K$, they base a direct summand of $\Omega_{R/k'} \otimes_R A$. Thus $\Omega_{A/k'}$ is also a
free $A$-module, and $\rank \Omega_{A/k'} + r = \rank \Omega_{R/k'} = \rank \Omega_{k/k'} + n$.

(3) $\Rightarrow$ (1) In the exact sequence ($\dagger$) the two terms
$\Omega_{R/k'} \otimes_R A$ and $\Omega_{A/k'}$ are both free modules, and the difference between
their ranks is $r$, so that $IR/I^2R$ maps onto a direct summand of
$\Omega_{R/k'} \otimes_R A$, which is a free module of rank $r$. Thus we can choose
$f_1, \ldots, f_r \in IR$ such that $df_1 \otimes 1, \ldots, df_r \otimes 1$ base this direct summand,
and then by NAK we see that $Rdf_1 + \cdots + Rdf_r$ is a direct summand of
$\Omega_{R/k'}$. Hence there exist $D_1, \ldots, D_r \in \Der_{k'}(R)$ such that
$\det(D_if_j) \notin M$. Therefore, by Theorem~\hyperref[thm:30.4]{4}, $A$ is regular.
\end{proof}

\begin{corollary*}
Let $k$ be a field and $S = k[X_1, \ldots, X_n]$; let $I$ be an ideal of $S$, and set
$B = S/I$. Define $U = \{p \in \Spec B \mid B_p \text{ is 0-smooth over } k\}$ and
$\Reg(B) = \{p \in \Spec B \mid B_p \text{ is regular}\}$. Then both $U$ and $\Reg(B)$ are open
subsets of $\Spec B$.
\end{corollary*}

\begin{proof}
Set $V' = \Spec B$, and let $V_1, \ldots, V_s$ be the irreducible components of $V'$. To
say that $P \in V_i \cap V_j$ for $i \neq j$ means just that $B_p$ has at least two minimal
prime ideals, and such points cannot belong to $\Reg(B)$, (nor a fortiori to $U$).
Thus first of all we can throw out the closed subset $W = \bigcup_{i \neq j}(V_i \cap V_j)$,
and therefore we need only prove that $V_i \cap U$ and $V_i \cap \Reg(B)$ are open in
$V_i$ for each $i$; we fix $i$, and set $\dim V_i = n - r$. Then by
Theorems~\hyperref[thm:30.3]{3} and \hyperref[thm:30.4]{4}, if we let $\Delta_1, \ldots, \Delta_s$ denote the images in
$B$ of the $r \times r$ minors of the Jacobian matrix $(\partial f_i/\partial X_j)$, where
$I = (f_1, \ldots, f_t)$, then $V_i - U$ is the intersection of $V_i$ with the closed subset of
$V'$ defined by the ideal $(\Delta_1, \ldots, \Delta_s)B$ of $B$, and is thus closed in $V_i$.
Using Theorem~\hyperref[thm:30.5]{5}, we can argue similarly for $\Reg(B)$; the only difference
is that, instead of one Jacobian matrix, we have to consider the closed subsets
of $V'$ defined by the ideal of $B$ generated by all the $r \times r$ minors of the
infinitely many Jacobian matrices $J(f_1, \ldots, f_t; D_\alpha, D_\beta, \ldots, D_\gamma,
\partial/\partial X_1, \ldots, \partial/\partial X_n)$, where $\{D_\alpha, \ldots, D_\gamma\}$ runs through all finite
subsets of the set of derivations $\{D_\gamma\}_{\gamma \in \Gamma}$ appearing in
Theorem~\hyperref[thm:30.5]{5}, (2).
\end{proof}

Theorems~\hyperref[thm:30.3]{3} and \hyperref[thm:30.5]{5} contain the result known as the \emph{Jacobian
criterion for regularity} in polynomial rings. There is some purpose in trying to
extend this to more general rings. In cases when the module of differentials is
not finitely generated, then the above method cannot be used as it stands, so we
approach the problem using modules of derivations.

Quite generally, let $A$ be an integral domain with field of fractions $L$; then
for an $A$-module $M$, we write $\rank_A M$ for the dimension over $L$ of the vector
space $M \otimes_A L$.

\begin{theorem}[\textbf{M.~Nomura}]\label{thm:30.6}
Let $(R, m)$ be an equicharacteristic $n$-dimensional regular local ring, and $R^*$
the completion of $R$; suppose that $k$ is a quasi-coefficient field of $R$ and $K$ a
coefficient field of $R^*$ containing $k$. Let $x_1, \ldots, x_n$ be a regular system of
parameters of $R$.
\begin{enumerate}[(i)]
  \item $R^* = K[[x_1, \ldots, x_n]]$, and if we write $\partial/\partial x_i$ for the partial
        derivatives in this representation, then $\Der_k(R^*) = \Der_k(R^*)$ is the free
        $R^*$-module with basis $\partial/\partial x_i$ for $1 \le i \le n$.
  \item The following conditions are equivalent:
        \begin{enumerate}[(1)]
          \item $\partial/\partial x_i$ maps $R$ into $R$ for $1 \le i \le n$, so that they can be
                considered as elements of $\Der_k(R)$;
          \item there exist $D_1, \ldots, D_n \in \Der_k(R)$ and $a_1, \ldots, a_n \in R$ such
                that $D_ia_j = \delta_{ij}$;
          \item there exist $D_1, \ldots, D_n \in \Der_k(R)$ and $a_1, \ldots, a_n \in R$ such
                that $\det(D_ia_j) \notin m$;
          \item $\Der_k(R)$ is a free $R$-module of rank $n$;
          \item $\rank \Der_k(R) = n$.
        \end{enumerate}
\end{enumerate}
\end{theorem}

\begin{proof}
(i) Since $K$ is 0-\'etale over $k$, any derivation of $R^*$ which vanishes on $k$ also
vanishes on $K$. If $D \in \Der_k(R^*)$, set $Dx_i = y_i$; then for any
$f \in R^* = K[[x_1, \ldots, x_n]]$ we have $D(f) = \sum_{i=1}^{n}(\partial f/\partial x_i) \cdot y_i$,
and hence $D = \sum y_i \partial/\partial x_i$. Conversely, for any given $y_i$ we can
construct a derivation by this formula, so that $\Der_k(R^*)$ is the free $R^*$-module
with basis $\partial/\partial x_1, \ldots, \partial/\partial x_n$.

(ii) (1) $\Rightarrow$ (2) $\Rightarrow$ (3) and (4) $\Rightarrow$ (5) are trivial. If (3) holds then
$D_1, \ldots, D_n$ are linearly independent over $R$ and over $R^*$, so that by (i), any
$D \in \Der_k(R)$ can be written as a combination $D = \sum c_iD_i$ of the $D_i$ with
coefficients in the field of fractions of $R^*$, but since $Da_j = \sum c_iD_i(a_j)$ for
$j = 1, \ldots, n$, we have $c_i \in R$; therefore $D_1, \ldots, D_n$ form a basis of
$\Der_k(R)$, which proves (4).

(5) $\Rightarrow$ (1) If $D_1, \ldots, D_n$ are linearly independent over $R$ then there exist
$a_1, \ldots, a_n \in R$ such that $\det(D_ia_j) \ne 0$. Therefore $D_1, \ldots, D_n$ are also
linearly independent over $R^*$. Thus writing $L^*$ for the field of fractions of
$R^*$ we can write $\partial/\partial x_i = \sum_j c_{ij}D_j$, with $c_{ij} \in L^*$. From this we
get $\delta_{ij} = \sum_\nu c_{i\nu}D_\nu x_j$, so that the matrix $(c_{ij})$ is the inverse of
$(D_jx_i)$, and $c_{ij} \in L$, where $L$ is the field of fractions of $R$; therefore
$(\partial/\partial x_i)(R) \subset L \cap R^* = R$.
\end{proof}

\begin{lemma}\label{lem:30.2}
Let $R$ be a regular ring and $P \in \Spec R$ with $\htop P = r$; then the following
two conditions are equivalent:
\begin{enumerate}[(1)]
  \item there exist $D_1, \ldots, D_r \in \Der(R)$ and $f_1, \ldots, f_r \in P$ such that
        $\det(D_if_j) \notin P$;
  \item for all $Q \in \Spec R$ with $\htop Q = s \le r$ such that $Q \subset P$ and
        $R_P/QR_P$ is regular, there exist $D_1, \ldots, D_s \in \Der(R)$ and
        $g_1, \ldots, g_s \in Q$ such that $\det(D_ig_j) \notin P$.
\end{enumerate}
\end{lemma}

\begin{proof}
(2) contains (1) as the special case $P = Q$; conversely, suppose that (1) holds.
By Theorem~\hyperref[thm:30.4]{4} we see that $f_1, \ldots, f_r$ is a regular system of parameters
of $R_P$. If $R_P/QR_P$ is regular then we can take $g_1, \ldots, g_s \in Q$ forming a
minimal basis of $QR_P$, and then take $g_{s+1}, \ldots, g_r \in P$ such that
$g_1, \ldots, g_r$ is a regular system of parameters of $R_P$. Then from
$\det(D_if_j) \notin P$ we deduce that $\det(D_ig_j) \notin P$, or in other words
$\rank J(g_1, \ldots, g_r; D_1, \ldots, D_r)(P) = r$. Therefore
\[
  \rank J(g_1, \ldots, g_s; D_1, \ldots, D_s)(P) = s.
\]
\end{proof}

If the above condition (1) holds, we say that the \emph{weak Jacobian condition} (WJ)
holds at $P$. If (WJ) holds at every $P \in \Spec R$ then we say that (WJ) holds in $R$.
If this holds, then for any $P$, $Q \in \Spec R$ with $Q \subset P$, setting $\htop Q = s$ we
have
\[
  R_P/QR_P \text{ is regular} \iff
  \text{there exist } D_1, \ldots, D_s \in \Der(R) \text{ and }
  f_1, \ldots, f_s \in Q \text{ such that } \det(D_if_j) \notin P,
\]
(the implication ($\Leftarrow$) is given by Theorem~\hyperref[thm:30.4]{4}). This statement is
the Jacobian criterion for regularity. We can use $\Der_k(R)$ in place of
$\Der(R)$, and we then write (WJ)$_k$.

\begin{theorem}\label{thm:30.7}
Let $(A, m)$ be an $n$-dimensional Noetherian local integral domain containing
$\mathbb{Q}$, and let $k \subset A$ be a subfield such that $\trdeg_k(A/m) = r < \infty$. Then
$\Der_k(A)$ is isomorphic to a submodule of $A^{n+r}$, and is therefore a finite
$A$-module, with
\[
  \rank \Der_k(A) \le \dim A + \trdeg_k(A/m).
\]
\end{theorem}

\begin{proof}
We write $A^*$ for the completion of $A$; choose a quasi-coefficient field $k'$ of $A$
containing $k$, and let $K$ be a coefficient field of $A^*$ containing $k'$. Let
$u_1, \ldots, u_r$ be a transcendence basis of $k'$ over $k$, and $x_1, \ldots, x_n$ a system
of parameters of $A$. We define $\varphi\colon \Der_k(A) \to A^{n+r}$ by
$\varphi(D) = (Du_1, \ldots, Du_r, Dx_1, \ldots, Dx_n)$; then $\varphi$ is $A$-linear, so that we
need only prove that it is injective. Suppose then that $Du_i = Dx_j = 0$ for all $i$
and $j$; then $D$ has a continuous extension to $A^*$, and this vanishes on
$B = K[[x_1, \ldots, x_n]]$. Now we do not know whether $A^*$ is an integral domain,
but it is finite as a $B$-module, so that any $a \in A$ is integral over $B$, and if
$f(X) \in B[X]$ has $a$ as a root, and has minimal degree, then $f(a) = 0$,
$f'(a) \ne 0$. Then $0 = D(f(a)) = f'(a) \cdot Da$, and $Da \in A$, so that, since a
non-zero element of $A$ cannot be a zero-divisor in $A^*$, we must have $Da = 0$.
Hence $D = 0$.
\end{proof}

\begin{remark}
If $k$ is an imperfect field then there are counter-examples even if
$k = A/m$: suppose $\chark k = p$ and $a \in k - k^p$, and set
$A = k[X, Y]_{(X, Y)}/(X^p + aY^p)$; then $\dim A = 1$ but
$\rank \Der_k(A) = 2$.
\end{remark}

\begin{theorem}\label{thm:30.8}
Let $(R, m)$ be a regular local ring containing $\mathbb{Q}$ and $k$ a quasi-coefficient
field of $R$. Then the following three conditions are equivalent:
\begin{enumerate}[(1)]
  \item (WJ)$_k$ holds at $m$;
  \item $\rank \Der_k(R) = \dim R$;
  \item (WJ)$_k$ holds at every $P \in \Spec R$.
\end{enumerate}
Furthermore, if these conditions hold, then for any $P \in \Spec R$, every element
of $\Der_k(R/P)$ is induced by an element of $\Der_k(R)$ and
\[
  \rank \Der_k(R/P) = \dim R/P.
\]
\end{theorem}

\begin{proof}
(1) $\Leftrightarrow$ (2) is known from Theorem~\hyperref[thm:30.6]{6}, and (1) is contained in
(3).

(1) $\Rightarrow$ (3) Write $R^*$ for the completion of $R$, and let $K$ be the coefficient
field of $R^*$ containing $k$; then by Theorem~\hyperref[thm:30.6]{6}, if
$x_1, \ldots, x_n$ is a regular system of parameters of $R$ then the derivations
$\partial/\partial x_i$ of $R^* = K[[x_1, \ldots, x_n]]$ belong to $\Der_k(R)$ for $1 \le i \le n$, and
form a basis of it. Now let $P \in \Spec R$, and write $\sigma\colon R \to R/P$ for the
natural homomorphism; to say that $D' \in \Der_k(R/P)$ is induced by
$D \in \Der_k(R)$ means that there is a commutative diagram
\[
\begin{tikzcd}
R \arrow[r, "\sigma"] \arrow[d, "D"'] & R/P \arrow[d, "D'"] \\
R \arrow[r, "\sigma"'] & R/P
\end{tikzcd}
\]
Suppose then that $D'$ is given; then $D'\sigma \in \Der_k(R, R/P)$, and this has a
unique extension to an element of $\Der_k(R^*, R^*/PR^*)$, so that it is uniquely
determined by its values on $x_1, \ldots, x_n$. Therefore if we choose
$b_1, \ldots, b_n \in R$ such that $D'(\sigma(x_i)) = \sigma(b_i)$, and set
$D = \sum b_i \partial/\partial x_i$, then $D'$ is induced by $D$.

Now $\Der_k(R, R/P)$ is a free $R/P$-module with basis $\bar{\sigma} \circ \partial/\partial x_i$ for
$1 \le i \le n$, and $\Der_k(R/P)$ can be identified with the submodule
\[
  N = \{D \in \Der_k(R, R/P) \mid D(f) = 0 \text{ for all } f \in P\}.
\]
Therefore, if $P = (f_1, \ldots, f_t)$ and $\htop P = r$ then
\[
  \rank \Der_k(R/P) = n - \rank J(f_1, \ldots, f_t; \partial/\partial x_1, \ldots,
      \partial/\partial x_n)(P);
\]
according to Theorem~\hyperref[thm:30.4]{4}, the right-hand side is $\ge n - r$, and by
Theorem~\hyperref[thm:30.7]{7} the left-hand side is $\le \dim R/P = n - r$, so we see that
\[
  \rank J(f_1, \ldots, f_t; \partial/\partial x_1, \ldots, \partial/\partial x_n)(P) = r,
\]
\[
  \rank \Der_k(R/P) = \dim R/P.
\]
\end{proof}

This theorem has many applications. For example, in the ring $R$ of convergent
power series in $n$ variables over $k = \mathbb{R}$ or $\mathbb{C}$ ($R$ is denoted by
$k\{X_1, \ldots, X_n\}$ in [N1], and elsewhere by $k\{X_1, \ldots, X_n\}$) we have
$\partial X_i/\partial X_j = \delta_{ij}$, so that the Jacobian criterion for regularity holds in
$R$. In characteristic $0$, the assumptions of the theorem are satisfied by the
formal power series ring over any field, or more generally by the formal power
series ring $R[[Y_1, \ldots, Y_n]]$ over a regular local ring $R$ which satisfies the
assumptions of the theorem. Even if $R$ is not local, but is a regular ring
containing a field $k$ of characteristic $0$, and such that the residue field at
every maximal ideal is algebraic over $k$, then if (WJ)$_k$ holds at every maximal
ideal of $R$, it in fact holds at every prime ideal of $R$, so that the Jacobian
criterion for regularity holds for example in rings such as
$k[X_1, \ldots, X_n] \otimes k[[Y_1, \ldots, Y_m]]$. A further extension of this theorem can be
found in Matsumura [2].

Now we are going to prove the results analogous to Theorem~\hyperref[thm:30.5]{5} for formal
power series rings over a field of characteristic $p$. This is a difficult theorem
obtained by Nagata [6]. First of all we have to do some preparatory work.

Let $k$ be a field and $k' \subset k$ a subfield; we say that $k'$ is \emph{cofinite} in $k$ if
$[k : k'] < \infty$. We say that a family $\mathcal{F} = \{k_\alpha\}_{\alpha \in I}$ of subfields of $k$ is
a \emph{directed family} if for any $\alpha, \beta \in I$ there exists $\gamma \in I$ such that
$k_\alpha, k_\beta \subset k_\gamma$.

Let $K$ be a field of characteristic $p$, and $k \subset K$ a subfield. Then there
exists a directed family $\mathcal{F} = \{k_\alpha\}_{\alpha \in I}$ of intermediate fields
$k \subset k_\alpha \subset K$ cofinite in $K$ such that $\bigcap_\alpha k_\alpha = k(K^p)$. To construct
this, we let $B$ be a fixed $p$-basis of $K$ over $k$, and let $I$ be the set of all
finite subsets of $B$; then we only need take $k_\alpha = k(K^p, B - \alpha)$ for
$\alpha \in I$.

\begin{lemma}\label{lem:30.3}
Let $K$ be a field, $(k_\alpha)_{\alpha \in I}$ a directed family of subfields of $K$, and set
$k = \bigcap_\alpha k_\alpha$. Then if $V$ is a vector space over $K$, and
$v_1, \ldots, v_n \in V$ are linearly independent over $k$, there exists $\alpha \in I$ such that
$v_1, \ldots, v_n$ are also linearly independent over $k_\alpha$.
\end{lemma}

\begin{proof}
For each $\alpha \in I$, write $q(\alpha)$ for the number of linearly independent
elements over $k_\alpha$ among $v_1, \ldots, v_n$; let $\alpha$ be such that $q(\alpha)$ is maximal,
and set $q = q(\alpha)$. Now if $q < n$, and we assume that $v_1, \ldots, v_q$ are linearly
independent over $k_\alpha$, we have $v_n = \sum_{i=1}^{q} c_iv_i$ with $c_i \in k_\alpha$. Since
the $v_i$ are linearly independent over $k$, at least one of the $c_i$ does not belong
to $k$, so that we can assume $c_1 \notin k$. Hence there exists $\beta \in I$ such that
$c_1 \notin k_\beta$, and also $\gamma \in I$ such that $k_\alpha, k_\beta \subset k_\gamma$, so that
$v_1, \ldots, v_q$ and $v_n$ are linearly independent over $k_\gamma$; this contradicts the
maximality of $q$, so that $q = n$.
\end{proof}

\begin{lemma}\label{lem:30.4}
Let $k \subset K$ be fields of characteristic $p$, and let $\mathcal{F} = \{k_\alpha\}_{\alpha \in I}$
be a directed family of intermediate fields $k \subset k_\alpha \subset K$; then the
following conditions are equivalent:
\begin{enumerate}[(1)]
  \item $\bigcap_\alpha k_\alpha(K^p) = k(K^p)$;
  \item the natural map $\Omega_{K/k} \to \varprojlim \Omega_{K/k_\alpha}$ is injective;
  \item if a finite subset $\{u_1, \ldots, u_n\} \subset K$ is $p$-independent over $k$, then
        it is also $p$-independent over $k_\alpha$ for some $\alpha$;
  \item there exists a $p$-basis $B$ of $K$ over $k$ such that every finite subset of
        $B$ is $p$-independent over $k_\alpha$ for some $\alpha$.
\end{enumerate}
\end{lemma}

\begin{proof}
(1) $\Rightarrow$ (3) If $\{u_1, \ldots, u_n\}$ is $p$-independent over $k$ then the $p^n$ monomials
$u_1^{i_1} \cdots u_n^{i_n}$ for $0 \le i_\nu < p$ are linearly independent over $k(K^p)$, so that by
the previous lemma they are also linearly independent over $k_\alpha(K^p)$ for some
$\alpha$.

(3) $\Rightarrow$ (4) is trivial; any $p$-basis will do.

(4) $\Rightarrow$ (2) If $0 \ne \omega \in \Omega_{K/k}$ and $B$ is a $p$-basis then there is a unique
expression $\omega = \sum c_id_{K/k}b_i$, with $\{b_1, \ldots, b_n\} \subset B$ a finite set and
$c_i \in K$. If we take $\alpha$ such that $b_1, \ldots, b_n$ are $p$-independent over
$k_\alpha$, then $d_{K/k_\alpha}b_i$ for $1 \le i \le n$ are linearly independent as elements of
$\Omega_{K/k_\alpha}$, so that $\omega$ has non-zero image in $\varprojlim \Omega_{K/k_\alpha}$.

(2) $\Rightarrow$ (1) Let $a \in K$ be such that $a \notin k(K^p)$; then $d_{K/k}a \ne 0$, so that
$d_{K/k}a \ne 0$ in $\Omega_{K/k_\alpha}$ for some $k_\alpha$, in other words $a \notin k_\alpha(K^p)$.
\end{proof}

\begin{lemma}\label{lem:30.5}
Let $k \subset K$ and $\mathcal{F} = \{k_\alpha\}_{\alpha \in I}$ be as in the previous lemma, and
assume that $\bigcap_\alpha k_\alpha(K^p) = k(K^p)$. Then if $L$ is an extension field of $K$
which is either separable over $K$ or finitely generated over $K$, we again have
$\bigcap_\alpha k_\alpha(L^p) = k(L^p)$.
\end{lemma}

\begin{proof}
(i) If $L/K$ is separable, choose a $p$-basis $B$ of $K/k$ and a $p$-basis $C$ of
$L/K$; then in view of the exact sequence (Theorem~\hyperref[thm:25.1]{25.1})
\[
  0 \to \Omega_{K/k} \otimes_K L \to \Omega_{L/k} \to \Omega_{L/K} \to 0,
\]
$B \cup C$ is a $p$-basis of $L/k$. If $b_1, \ldots, b_m \in B$ and $c_1, \ldots, c_n \in C$ are
finitely many distinct elements, then by the previous lemma, $b_1, \ldots, b_m$ are
$p$-independent over some $k_\alpha$, and from this (replacing $k$ by $k_\alpha$ in the
above exact sequence) we see that $b_1, \ldots, b_m, c_1, \ldots, c_n$ are
$p$-independent over $k_\alpha$.

(ii) If $L/K$ is finitely generated, then since any finitely generated extension
can be obtained as a succession of elementary extensions of type
\begin{enumerate}[(a)]
  \item $L = K(x)$ with $x$ separable over $K$,
  or
  \item $L = K(x)$ with $x^p = a \in K$,
\end{enumerate}
we need only consider case (b). We further divide this into two subcases: in
the first, $d_{K/k}a = 0$, and then from Theorem~\hyperref[thm:25.2]{25.2} and the fact that
$L \cong K[X]/(X^p - a)$ we get $\Omega_{L/k} = (\Omega_{K/k} \otimes L) \oplus Ldx$; from this
one sees easily that $\Omega_{L/k} \to \varprojlim \Omega_{L/k_\alpha}$ is injective. In the second
subcase, $d_{K/k}a \ne 0$ and now we have
$\Omega_{L/k} \cong ((\Omega_{K/k} \otimes L)/L \cdot d_{K/k}a) \oplus Ldx$. Hence if we take a
$p$-basis of $K/k$ of the form $\{a\} \cup B'$ with $a \notin B'$, then $\{x\} \cup B'$ is a
$p$-basis of $L/k$. Now $L/k$ satisfies condition (4) of the previous lemma, since
for $b_1, \ldots, b_m \in B'$, if we choose $\alpha$ such that $\{a, b_1, \ldots, b_m\} \subset K$
is $p$-independent over $k_\alpha$, then $\{x, b_1, \ldots, b_m\} \subset L$ is $p$-independent
over $k_\alpha$.
\end{proof}

\begin{lemma}\label{lem:30.6}
Let $K$ be a field of characteristic $p$, and let $\{K_\alpha\}$ be a directed family of
cofinite subfields $K_\alpha \subset K$ such that $\bigcap_\alpha K_\alpha = K^p$. Then if $L$ is a
finite field extension of $K$, there exists $\alpha$ such that
\[
  \rank_L \Omega_{L/K_\alpha} = \rank_L \Omega_{L/K}
\]
for all subfields $K' \subset K_\alpha$ with $[K : K'] < \infty$.
\end{lemma}

\begin{proof}
Let $K = K_0 \subset K_1 \subset \cdots \subset K_t = L$ be a chain of intermediate
fields such that $K_i = K_{i-1}(x_i)$, with $x_i$ either separable algebraic over
$K_{i-1}$ or $x_i^p \in K_{i-1}$. Then by the previous lemma we have
$\bigcap_\alpha K_\alpha(L^p) = L^p$, so that we need only prove the lemma for $t = 1$;
hence suppose that $L = K(x)$. If $L$ is separable over $K$ then in view of
$\Omega_{L/K'} = \Omega_{K/K'} \otimes_K L$, the assertion is clear (any $\alpha$ will do). Thus
suppose that $x^p = a \in K$, but $a \notin K^p$; then there exists $\alpha$ such that
$a \notin K_\alpha$. If $K' \subset K_\alpha$, then computing by means of
$L = K[X]/(X^p - a)$ we see that
\[
  \Omega_{L/K'} = ((\Omega_{K/K'} \otimes_K L) \oplus Ldx)/L \cdot d_{K/K'}a,
\]
and if $\rank \Omega_{L/K'} < \infty$ then $\rank \Omega_{L/K'} = \rank \Omega_{L/K}$.
\end{proof}

\begin{theorem}[\textbf{Nagata}]\label{thm:30.9}
Let $k$ be a field of characteristic $p$, $S = k[Y_1, \ldots, Y_m]$ and $P \in \Spec S$;
suppose that $\mathcal{F} = \{k_\alpha\}_{\alpha \in I}$ is a directed family of cofinite subfields of
$k$ such that $\bigcap_\alpha k_\alpha = k^p$. Then there exists $\alpha \in I$ such that for
every intermediate field $k^p \subset k' \subset k_\alpha$ with $[k : k'] < \infty$ the
following formula holds:
\[
  \rank \Der_{k'}(S/P) = \dim(S/P) + \rank \Der_{k'}(k).
\]
\end{theorem}

\begin{proof}
Set $A = S/P$, let $L$ be the field of fractions of $A$, and $\dim A = n$. Choose a
system of parameters $x_1, \ldots, x_n$ of $A$, set $B = k[x_1, \ldots, x_n]$, let $K$ be the
field of fractions of $B$ and $m_B$ its maximal ideal. Then $A$ is a finite $B$-module,
and hence $[L : K] < \infty$. If $k'$ is an intermediate field $k^p \subset k' \subset k$
with $[k : k'] = p^\rho < \infty$, and if $u_1, \ldots, u_\rho$ is a $p$-basis of $k$ over $k'$,
then $\{u_1, \ldots, u_\rho, x_1, \ldots, x_n\}$ is a $p$-basis of $B$ over
$C' = k'[x_1^p, \ldots, x_n^p]$, in the sense that $B$ is the free $C'$-module with basis
the set of $p$-monomials in $u_1, \ldots, u_\rho, x_1, \ldots, x_n$. A derivation from $B$ to
$B$ is continuous in the $m_B$-adic topology, and any element of $\Der_{k'}(B)$ is $0$ on
$C'$. Therefore $\Der_{k'}(B) = \Der_{k'}(B) = \Hom_B(\Omega_{B/C'}, B)$, and $\Omega_{B/C'}$ is
the free $B$-module of rank $n + \rho$ with basis $du_1, \ldots, du_\rho, dx_1, \ldots, dx_n$.
Therefore $\Der_{k'}(B)$ is also a free $B$-module of rank $n + \rho$, and the theorem
holds in case $A = B$.

We write $F'$ for the field of fractions of $C'$, and for $k_\alpha \in \mathcal{F}$ we set
$C_\alpha = k_\alpha[x_1^p, \ldots, x_n^p]$ and write $K_\alpha$ for the field of fractions of
$C_\alpha$. As above we have
\[
  \Der_{k'}(A) = \Der_{k'}(A) = \Hom_A(\Omega_{A/C'}, A),
\]
and since $A$ is a finite $C'$-module, $\Omega_{A/C'}$ is a finite $A$-module. Hence
\[
  \Der_{k'}(A) \otimes_A L = \Hom_L(\Omega_{A/C'} \otimes_A L, L) = \Hom_L(\Omega_{L/F'}, L),
\]
so that $\rank_L \Der_{k'}(A) = \rank_L \Omega_{L/F'}$. Moreover,
\[
  n + \rank \Der_{k'}(k) = \rank \Omega_{L/F'} = \rank \Omega_{K/F'},
\]
so that the conclusion of the theorem can be rewritten
\[
  \rank \Omega_{L/F'} = \rank \Omega_{K/F'}.
\]
Any element of $K_\alpha$ can be written with its denominator in
$k^p[x_1^p, \ldots, x_n^p] = B^p$. However, $B$ is faithfully flat over $C_\alpha$, so that
$K_\alpha \cap B = C_\alpha$. From this one deduces easily that
$\bigcap_\alpha K_\alpha = K^p$. Thus the theorem follows from the previous lemma.
\end{proof}

\begin{theorem}[\textbf{Nagata's Jacobian criterion}]\label{thm:30.10}
Let $k$ be a field of characteristic $p$, $S = k[[X_1, \ldots, X_n]]$, and let $I$, $P$ be
ideals of $S$ such that $I \subset P \in \Spec S$; set $S_P = R$, $R/IR = A$, and suppose
that $\htop IR = r$ and $I = (f_1, \ldots, f_t)$. Choose a $p$-basis $\{u_\gamma\}_{\gamma \in \Gamma}$ of
$k$, and define $D_\gamma \in \Der(k)$ by $D_\gamma(u_{\gamma'}) = \delta_{\gamma\gamma'}$. We make
any element $D \in \Der(k)$ act on the coefficients of power series, thus
extending $D$ to an element of $\Der(S)$.

Then the following conditions are equivalent:
\begin{enumerate}[(1)]
  \item $A$ is a regular local ring;
  \item there exists a finite number of elements $\alpha, \beta, \ldots, \gamma \in \Gamma$
        such that
        \[
          \rank J(f_1, \ldots, f_t; D_\alpha, \ldots, D_\gamma,
              \partial/\partial X_1, \ldots, \partial/\partial X_n)(P) = r.
        \]
\end{enumerate}
\end{theorem}

\begin{proof}
(2) $\Rightarrow$ (1) follows from Theorem~\hyperref[thm:30.4]{4}.

(1) $\Rightarrow$ (2) As we see from the proof of Lemma~\hyperref[lem:30.2]{2}, we need only
prove (2) in the case $I = P$. We let $\mathcal{F}$ be the family of subfields of $k$
obtained by adjoining all but a finite number of $\{u_\gamma\}_{\gamma \in \Gamma}$ to $k^p$;
then the conditions of Theorem~\hyperref[thm:30.9]{9} are satisfied, and there exists
$k' \in \mathcal{F}$ such that
\[
  \rank \Der_{k'}(S/P) = n - \htop P + \rank \Der_{k'}(k).
\]
If we set $[k : k'] = p^s$ then by construction there exist
$\gamma_1, \ldots, \gamma_s \in \Gamma$ such that $k = k'(u_{\gamma_1}, \ldots, u_{\gamma_s})$; now set
$C = k'[[X_1, \ldots, X_n]]$, so that $\Omega_{S/C}$ is the free $S$-module with basis
$du_{\gamma_1}, \ldots, du_{\gamma_s}, dX_1, \ldots, dX_n$, and arguing as in the proof of
Theorem~\hyperref[thm:30.8]{8}, we see that the derivations of $S/P$ over $k'$ are all induced by
derivations of $S$ over $k'$, and that if $g_1, \ldots, g_m$ are generators of $P$, we have
\[
  \rank J(g_1, \ldots, g_m; D_{\gamma_1}, \ldots, D_{\gamma_s},
      \partial/\partial X_1, \ldots, \partial/\partial X_n)(P) = \htop P.
\]
\end{proof}

\begin{remark}
Conditions (1) and (2) in Theorem~\hyperref[thm:30.10]{10} are of the same form as the
corresponding conditions in Theorem~\hyperref[thm:30.5]{5}. Condition (3) of
Theorem~\hyperref[thm:30.5]{5} is not applicable as it stands to the present situation, since
$\Omega_{A/k}$ is in general not a finite $A$-module. However, in general for a module
$M$ over a local ring $(R, m)$, if we write (in temporary notation)
$\tilde{M} = M/\bigcap_\nu m^\nu M$ for the associated separated module, then for any
$R$-module $N$ which is separated ($N = \tilde{N}$) we have
$\Hom_R(\tilde{M}, N) = \Hom_R(M, N)$. Therefore in the situation of the above theorem
we have $\Der_k(S) = \Hom_S(\tilde{\Omega}_{S/k}, S)$ and $\Der_k(A) = \Hom_A(\tilde{\Omega}_{A/k}, A)$,
and moreover $\tilde{\Omega}_{S/k}$ is a free $S$-module of rank $\rho + \rank \Omega_{k/\pi}$.
From this, using the same argument as in the proof of Theorem~\hyperref[thm:30.5]{5}, we see
that replacing $\Omega_{A/k}$ by $\tilde{\Omega}_{A/k}$ in (3) of Theorem~\hyperref[thm:30.5]{5}, this
condition is equivalent to (1) and (2) of Theorem~\hyperref[thm:30.10]{10}. Verifying this is a
suitable exercise for the reader.
\end{remark}

\begin{corollary*}
Let $A$ be a complete Noetherian local ring; then $\Reg(A)$ is an open subset of
$\Spec A$.
\end{corollary*}

\begin{proof}
If $A$ is equicharacteristic, and $k$ is a coefficient field of $A$, then $A$ is of the
form $A = S/I$ with $S = k[[X_1, \ldots, X_n]]$, so that by Theorems~\hyperref[thm:30.8]{8} and
\hyperref[thm:30.10]{10}, we see as in the corollary of Theorem~\hyperref[thm:30.5]{5} that $\Reg(A)$ is open.

If $A$ is of unequal characteristic, then by Theorem~\hyperref[thm:24.4]{24.4}, it is enough to
prove, under the assumption that $A$ is an integral domain, that $\Reg(A)$
contains a non-empty open subset of $\Spec A$. Now by
Theorem~\hyperref[thm:29.5]{29.5}, $A$ contains a regular local ring $B$, and is a finite module
over $B$, so that if we let $L$ and $K$ be the fields of fractions of $A$ and $B$,
then $L$ is a finite extension of $K$, and is separable since $\chark K = 0$. Replacing
$A$ by $A_b$, and $B$ by $B_b$ for some suitable $0 \ne b \in B$ we can assume that $A$ is
a free $B$-module (although $A$ and $B$ are no longer local rings, $B$ remains
regular). Suppose that $\omega_1, \ldots, \omega_m$ are a basis of $A$ as a $B$-module, and
consider the discriminant
\[
  d = \det(\tr_{L/K}(\omega_i\omega_j)).
\]
Let us prove that if $P \in \Spec A$ is such that $d \notin P$ then $P \in \Reg(A)$. Set
$p = P \cap B$; then $A_p$ is flat over $B_p$, and $B_p$ is regular, so that we need only
prove that the fibre $A_p \otimes_B \kappa(p)$ is regular. Now $A = \sum B\omega_i$, so
that $A \otimes_B \kappa(p) = \sum \kappa(p)\bar{\omega}_i$, and in $\kappa(p)$ we have
$\det(\tr(\bar{\omega}_i\bar{\omega}_j)) = \bar{d} \ne 0$; hence $A \otimes_B \kappa(p)$ is reduced,
and is therefore a direct product of fields. Therefore $A_p \otimes_B \kappa(p)$ is a
field; this proves that $\Reg(A)$ contains a non-empty open subset of $\Spec A$.
\end{proof}

% ============ 习题 ============
\section*{Exercises to \S\ref{sec:30}}

\begin{exercise}\label{ex:30.1-ex}
Let $(A, m)$ be a complete Noetherian local ring, and $D = (D_0, D_1, \ldots) \in
\mathrm{HS}(A)$. Suppose that $x \in m$ satisfies $D_1x = 1$ and $D_ix = 0$ for $i > 0$,
that is $E_t(x) = x + t$, and define $\varphi = E_{-t}$ by
$\varphi(a) = \sum_i (-x)^iD_ia$. Then $\varphi$ is an endomorphism of $A$ to $A$, with
$\Ker \varphi = xA$; and if we set $C = \Im \varphi$ then $A = C[x] \cong C[[X]]$.
\end{exercise}

\begin{exercise}\label{ex:30.2-ex}
If the conditions of Theorem~\hyperref[thm:30.5]{5} hold, is it true that $A$ is 0-smooth over
the $k'$ appearing in (3)?
\end{exercise}

\begin{exercise}\label{ex:30.3-ex}
Are conditions (1) and (3) of Theorem~\hyperref[thm:30.3]{3} still equivalent if
$S = k[[X_1, \ldots, X_n]]$?
\end{exercise}

\begin{exercise}\label{ex:30.4-ex}
Let $R$ be a regular ring containing a field of characteristic $0$; if (WJ) holds in
$R$ then it also holds in the polynomial ring $R[X_1, \ldots, X_n]$.
\end{exercise}

